% AUTOGENERATED DRAFT from manuscript source: manuscript/chapters/03a-warpx-initialization.md
% Converter: scripts/convert_markdown_chapter_to_tex.py
% Review gate: docs/latex-migration-plan.md (Phase 2 per-chapter gate)

\chapter{3A. WarpX 初始化链：从 \texttt{InitData()} 到初始粒子和外部场}
\label{chap:03a-warpx-initialization}

本章把 \cpath{WarpX::InitData()} 展开成一条完整的初始化链。它补足第 3 章中“初始化”只作为主循环前置步骤的不足：这里开始逐块解释 fresh run / restart、AMR level 初始化、外部场、species 注入器、粒子创建 kernel、Gaussian beam、openPMD 文件注入和 projection divergence cleaning。

阅读初始化代码时，应先围绕四个主入口建立因果关系：\cpath{InitData()} 决定 fresh/restart 分叉，\cpath{InitFromScratch()} 建立初始 level，\cpath{InitDiagnostics()} 准备可观察输出，\cpath{AddExternalFields()} 把外部场加入初态。使用不同 WarpX 版本时，优先按这些函数的职责和调用关系检索，不要把行号或分支名称当作初始化语义。

% <-- table block (source: 03a-warpx-initialization L7)
\begin{tabularx}{\textwidth}{LL}
\toprule
读者问题 & 首先追踪的对象 \\
\midrule
\endfirsthead
\toprule
读者问题 & 首先追踪的对象 \\
\midrule
\endhead
程序启动后哪些参数在构造对象前已被锁定？ & 参数读取、\cpath{WarpX::WarpX()} 与 \cpath{ReadParameters()} \\
参数、level 字段和 \cpath{InitData()} 的顺序如何落到对象上？ & \cpath{InitFromScratch()}、level allocation 与 field data \\
外场、species、粒子创建与散度修正分别如何进入初态？ & \cpath{AddExternalFields()}、\cpath{PlasmaInjector}、粒子创建与 projection cleaner \\
哪些案例能区分不同初始化路径？ & initial distribution、space charge、external field 与 restart 的观察量 \\
\bottomrule
\end{tabularx}

\subsection{阅读路线：先构造初态，再追踪分支}

初始化代码同时涉及参数、AMR、粒子、外部场和 diagnostics。第一次阅读应先建立“第一步推进之前哪些状态必须已经存在”的主线，再进入特定 injector 或 I/O 分支：

\begin{enumerate}
\item \textbf{先读 3A.1--3A.3。} 区分构造期、bootstrap、fresh run 与 restart，写清哪些参数会在对象构造或 \cpath{InitData()} 分叉前锁定；这一步回答初态从哪一个全局配置开始。
\item \textbf{再读 3A.4--3A.5。} 沿 level allocation、field data、PML、diagnostics 与外部场建立网格侧初态，特别区分叠加到网格的 external field 和在 gather 时供粒子消费的 external field；这一步回答第一步能读到哪些场对象。
\item \textbf{按输入类型读 3A.6--3A.12。} \cpath{PlasmaInjector}、体注入、Gaussian beam、openPMD 和 projection cleaner 是不同的初态构造路径。选择其中一条时，先记录它创建的粒子/场对象、时间位置和限制，而不是把它们统称为“加载初始条件”。
\item \textbf{最后读 3A.13--3A.16。} 用匹配的 regression、历史最小骨架和练习检验初态是否真的进入 \cpath{Evolve()}；结论必须区分输入/接口覆盖、writer 输出和物理 observable。
\end{enumerate}

这条路线的停止条件是：读者能够从一个输入参数追踪到它创建的初态对象，并说明第一个时间步的哪个阶段会消费该对象。

\section{3A.1 初始化链为什么值得单独成章}

理解初始化时，先把以下三层边界分开：

\begin{enumerate}
\item \cpath{WarpX::WarpX()} 构造期只完成参数读取和跨 level 外壳创建。
\item \cpath{InitData()} 才在 fresh run / restart 之间分叉，并把 AMR level、粒子、诊断、PML、外场和初始静电/磁静场组织到第一步之前。
\item species 注入、外部场、projection cleaning、openPMD/Gaussian beam 分别有自己的参数、时间层和验证边界，不能笼统地称为“初始化完成”。
\end{enumerate}

PIC 程序的时间推进方程通常写成：

\[
f^n,\mathbf{E}^n,\mathbf{B}^n
\longrightarrow
f^{n+1},\mathbf{E}^{n+1},\mathbf{B}^{n+1}.
\]

但第一个时间步之前必须先构造一个满足几何、边界、粒子权重、外部场、离散约束和并行布局的初态。WarpX 的初始化不是“读入参数然后开始跑”，而是完成以下任务：

\begin{enumerate}
\item 选择 fresh run 或 checkpoint restart。
\item 建立 AMReX level、field MultiFab、PML、EB 和 solver 数据结构。
\item 初始化 diagnostics 和 reduced diagnostics。
\item 创建 species、解析密度/动量/位置分布，并生成宏粒子。
\item 读入或解析外部场，并把外部场叠加到 self field 或 particle external field。
\item 对外部 \cpath{A/B} 场做 projection divergence cleaning。
\item 输出第 0 步 diagnostics，并检查 guard cell、solver 配置和 known issue。
\end{enumerate}

本章先建立从程序启动到第一个时间步的主干；需要实现级细节时，再沿源码文件、函数和调用关系逐项核对。

\section{3A.2 启动层先于 \texttt{InitData()}：MPI、AMReX、FFT、PETSc 与启动前提}

许多初始化说明从 \cpath{WarpX::InitData()} 往后讲，但 WarpX 真正的初始化链更早就开始了。\cpath{main.cpp} 最外层先做的是：

% <-- fenced code (cpp)
\begin{codeblock}
int
main (int argc, char* argv[]) {
    warpx::initialization::initialize_external_libraries(argc, argv);
    {
        auto& warpx = WarpX::GetInstance();
        warpx.InitData();
        warpx.Evolve();
        WarpX::Finalize();
    }
    warpx::initialization::finalize_external_libraries();
}
\end{codeblock}

这说明 \cpath{InitData()} 之前还有一层独立的 bootstrap 逻辑，负责：

\begin{enumerate}
\item 初始化 MPI、AMReX、FFT，以及可选 PETSc。
\item 覆盖一批 AMReX 默认 parser 行为。
\item 解析 geometry、periodicity 和整数数组表达式。
\item 锁定 \cpath{geometry.dims}、moving window 和 warning policy 这类全局运行前提。
\end{enumerate}

这一层主要分布在：

\begin{itemize}
\item \cpath{Initialization/WarpXAMReXInit.*}
\item \cpath{Initialization/WarpXInit.*}
\item \cpath{WarpX::MakeWarpX()}
\end{itemize}

\subsection{3A.2.1 参数系统先分命名空间，再分读取语义}

WarpX 的参数不是“一个大表一次性读完”。它先按 \cpath{ParmParse} 前缀拆成局部命名空间，然后再决定是：

\begin{itemize}
\item 立即求值成数值；
\item 保留 parser 字符串，稍后编译；
\item 解释成 interval 切片；
\item 还是先解析成枚举型算法选择。
\end{itemize}

最上游的运行控制参数就是无前缀直接读取：

% <-- fenced code (cpp)
\begin{codeblock}
const ParmParse pp;// Traditionally, max_step and stop_time do not have prefix.
utils::parser::queryWithParser(pp, "max_step", max_step);
utils::parser::queryWithParser(pp, "stop_time", stop_time);
\end{codeblock}

而大部分初始化参数则先按前缀取视图：

% <-- fenced code (cpp)
\begin{codeblock}
const ParmParse pp_amr("amr");
const ParmParse pp_algo("algo");
ParmParse const pp_warpx("warpx");
const ParmParse pp_geometry("geometry");
const ParmParse pp_boundary("boundary");
const ParmParse pp_psatd("psatd");
\end{codeblock}

接下来又会分成几种读取壳：

\begin{enumerate}
\item \cpath{queryWithParser/queryArrWithParser}
\begin{itemize}
\item 适合 \cpath{warpx.cfl}、\cpath{const\_dt}、\cpath{amr.n\_cell} 这类“读入就该变成数值”的参数；
\end{itemize}
\item \code{Store_parserString + makeParser}
\begin{itemize}
\item 适合 \code{ref_patch_function(x,y,z)}、外场 parser、diagnostics 过滤函数这类“先保留表达式，再在别处编译执行”的参数；
\end{itemize}
\item \cpath{IntervalsParser}
\begin{itemize}
\item 适合 \cpath{sort\_intervals}、\cpath{dt\_update\_interval}、diagnostics \cpath{intervals} 这类切片语法，不应误看成普通整数；
\end{itemize}
\item \code{query_enum_sloppy + AMREX_ENUM}
\begin{itemize}
\item 适合 \cpath{algo.evolve\_scheme}、\cpath{warpx.do\_electrostatic}、\cpath{algo.current\_deposition} 这类算法枚举入口。
\end{itemize}
\end{enumerate}

这层结构解释了为什么同样都出现在 \cpath{parameters.rst} 里，\cpath{max\_step}、\code{ref_patch_function(x,y,z)} 和 \cpath{algo.evolve\_scheme} 在源码中会走三种完全不同的消费链。后续参数索引如果想真正有用，就必须至少区分：

\begin{itemize}
\item 参数前缀；
\item 第一次被哪个 \cpath{ParmParse} 命名空间读取；
\item 它属于数值参数、parser 字符串、interval 还是枚举型算法选择。
\end{itemize}

例如 \cpath{WarpXAMReXInit.cpp} 并不是简单调用 \cpath{amrex::Initialize()}，而是把一个缺省值覆盖回调一起传进去：

% <-- fenced code (cpp)
\begin{codeblock}
amrex::Initialize(
    argc,
    argv,
    build_parm_parse,
    MPI_COMM_WORLD,
    ::overwrite_amrex_parser_defaults
);
\end{codeblock}

这个回调会统一注入常量、改写 \cpath{amrex.abort\_on\_out\_of\_gpu\_memory}、\cpath{amrex.the\_arena\_is\_managed}、\cpath{amrex.omp\_threads}、粒子 tiling 缺省值，以及由 \cpath{boundary.field\_*} / \cpath{boundary.particle\_*} 反推 \cpath{geometry.is\_periodic}。换句话说，WarpX 从程序一启动就已经不是“AMReX 原生默认设置”。

同一个文件还会在 AMReX 初始化后立即把 \cpath{geometry.prob\_lo/prob\_hi}、\cpath{amr.n\_cell}、\cpath{max\_grid\_size*}、\cpath{blocking\_factor*} 这类可能带表达式的输入预解析并写回 parser，保证后面的 \cpath{Geometry}、\cpath{warpx\_job\_info} 和 \cpath{yt} 读到的是数值结果，而不是未展开的表达式字符串。

\subsection{3A.2.2 文档 alias、AMReX-owned 参数与 WarpX 自有 parser 不是一回事}

参数索引继续往下清理后，会发现还有一类参数不能简单用“有没有 grep 到同名字符串”来理解。典型例子是：

\begin{itemize}
\item \cpath{geometry.prob\_lo/hi}
\item \cpath{boundary.field\_lo/hi}
\item \cpath{boundary.potential\_lo/hi\_x/y/z}
\item \cpath{psatd.nox/noy/noz}
\item \cpath{qed\_schwinger.xmin/ymin/zmin/xmax/ymax/zmax}
\end{itemize}

这些名字在文档里是 grouped alias，但源码里真实读取的仍然是拆开的 key。以 \cpath{geometry.prob\_lo/hi} 为例，WarpX 真正做的是：

% <-- fenced code (cpp)
\begin{codeblock}
utils::parser::getArrWithParser(
    pp_geometry, "prob_lo", prob_lo, 0, AMREX_SPACEDIM);
utils::parser::getArrWithParser(
    pp_geometry, "prob_hi", prob_hi, 0, AMREX_SPACEDIM);

pp_geometry.addarr("prob_lo", prob_lo);
pp_geometry.addarr("prob_hi", prob_hi);
\end{codeblock}

所以 \cpath{geometry.prob\_lo/hi} 只是文档层“成对参数”的写法，真正 parser key 仍然是 \cpath{prob\_lo} 和 \cpath{prob\_hi}。\cpath{boundary.field\_lo/hi}、\cpath{boundary.particle\_lo/hi}、\cpath{boundary.potential\_lo/hi\_x/y/z}、\cpath{psatd.nox/noy/noz} 和 Schwinger 区域边界框也都属于这一类。

还要再区分另一类：参数确实出现在 WarpX 手册里，但 WarpX 并不直接读取，而是由 AMReX 自己消费，WarpX 只消费结果。例如 \cpath{amr.ref\_ratio} / \cpath{amr.ref\_ratio\_vect} 和 \cpath{amrex.async\_out} / \cpath{amrex.async\_out\_nfiles}。本书引用的 WarpX 源树没有独立的：

% <-- fenced code (cpp)
\begin{codeblock}
ParmParse("amr").query("ref_ratio", ...)
ParmParse("amrex").query("async_out", ...)
\end{codeblock}

但后续代码会直接消费已经构造好的 \cpath{ref\_ratios}，或者在 plotfile I/O 邻近层继续设置 WarpX 自己的 \cpath{field\_io\_nfiles} / \cpath{particle\_io\_nfiles}。因此，参数索引如果要稳定，至少要区分三类：

\begin{enumerate}
\item WarpX 自身直接 parse；
\item WarpX 子对象 parse；
\item AMReX-owned 输入，WarpX 只消费 materialized 结果。
\end{enumerate}

另一条关键链在 \cpath{WarpX::MakeWarpX()}：

% <-- fenced code (cpp)
\begin{codeblock}
warpx::initialization::check_dims();
warpx::initialization::read_moving_window_parameters(...);
ConvertLabParamsToBoost();
parse_field_boundaries();
parse_particle_boundaries(...);
CheckGriddingForRZSpectral();
m_instance = new WarpX();
\end{codeblock}

因此单例构造前就已经锁定了四类全局事实：

\begin{itemize}
\item 当前可执行文件与 \cpath{geometry.dims} 是否匹配；
\item moving window 是否开启、方向和速度是什么；
\item field / particle boundary 类型是什么；
\item RZ spectral 的 gridding 约束是否满足。
\end{itemize}

例如 \cpath{check\_dims()} 直接把编译维度和 \cpath{geometry.dims} 做强一致断言，而 \cpath{read\_moving\_window\_parameters()} 会把输入文件里的 \cpath{moving\_window\_v} 从“以光速为单位的无量纲参数”转换成真正的物理速度 \codeesc{v = (\textbackslash\{\}cdots)c}，再写进 \cpath{WarpX} 的全局静态状态。

这里还要补一层工程来源。\cpath{check\_dims()} 能做这种强断言，前提不是“WarpX 在运行时随意切换几何”，而是构建系统已经先把几何变体锁死了。当前主构建链在顶层 \cpath{CMakeLists.txt} 中通过 \cpath{WarpX\_DIMS} 生成一组按维度后缀分裂的 \codeesc{lib\_\$\{SD\}} / \codeesc{pyWarpX\_\$\{SD\}} target；旧 GNUmake 链则通过 \cpath{DIM}、\cpath{USE\_RZ} 以及 \cpath{-DWARPX\_DIM\_3D}、\cpath{-DWARPX\_DIM\_XZ}、\cpath{-DWARPX\_DIM\_RZ} 这组宏编码几何变体。因此初始化阶段读到的 \cpath{geometry.dims} 实际是在和“当前可执行文件已经按哪一种几何编译出来”做一致性检查，而不是在决定后面要不要临时切换到另一种维度。

再往下，\cpath{WarpX} 构造函数开头还会调用：

% <-- fenced code (cpp)
\begin{codeblock}
warpx::initialization::initialize_warning_manager();
\end{codeblock}

也就是在任何 \cpath{ReadParameters()} 和 \cpath{InitData()} 之前，先读入：

\begin{itemize}
\item \cpath{warpx.always\_warn\_immediately}
\item \cpath{warpx.abort\_on\_warning\_threshold}
\end{itemize}

这说明 warning manager 在 WarpX 里也是初始化启动层的一部分，而不是后面运行时再临时决定的日志选项。

这一层再往下细分，还要注意 \cpath{MakeWarpX()} 和构造期 \cpath{ReadParameters()} 之间并不是简单前后顺序，而是“前者已经开始改写后者将要读取的参数”。例如 \cpath{ConvertLabParamsToBoost()} 不只是保存 \cpath{gamma\_boost}，而是会先把：

\begin{itemize}
\item \cpath{geometry.prob\_lo/prob\_hi}
\item \cpath{warpx.fine\_tag\_lo/fine\_tag\_hi}
\item \cpath{slice.dom\_lo/dom\_hi}
\end{itemize}

按 boosted-frame 规则直接回写进 parser。若同时启用 moving window，它计算的转换系数还会显式依赖 \cpath{moving\_window\_v/c}。因此 boosted-frame 与 moving-window 的第一次耦合，其实发生在 \cpath{ReadParameters()} 之前，而不是发生在后面的 \cpath{Evolve()}。

同样，\cpath{CheckGriddingForRZSpectral()} 也不是单纯做断言。对 \code{WARPX_DIM_RZ + algo.maxwell_solver = psatd} 这条组合，它会在构造 \cpath{WarpX} 对象之前直接改写：

\begin{itemize}
\item \cpath{amr.blocking\_factor\_x/y}
\item \cpath{amr.max\_grid\_size\_x/y}
\end{itemize}

并要求 longitudinal 方向至少满足“每个 MPI rank 至少有一个 block，且每个 rank 至少有 8 个 z cells”的约束。也就是说，构造函数里的 \cpath{ReadParameters()} 读到的并不总是用户输入文件里的原始 AMR 分块值，而往往已经是启动层重写过的 spectral-compatible 版本。

把这一点看清之后，初始化启动层就能更准确地拆成两段：

\begin{enumerate}
\item 预初始化与 parser 改写段：
\begin{itemize}
\item \cpath{initialize\_external\_libraries()}
\item \cpath{amrex\_init()}
\item \cpath{parse\_geometry\_input()}
\item \cpath{read\_moving\_window\_parameters()}
\item \cpath{ConvertLabParamsToBoost()}
\item \cpath{CheckGriddingForRZSpectral()}
\end{itemize}
\item 构造与运行态落地段：
\begin{itemize}
\item \cpath{initialize\_warning\_manager()}
\item \cpath{ReadParameters()}
\item \cpath{BackwardCompatibility()}
\item \cpath{InitEB()}
\item \cpath{MultiParticleContainer} / \cpath{ParticleBoundaryBuffer} / solver objects` 的真正创建
\end{itemize}
\end{enumerate}

这个分层很有用，因为后面再读 boosted-frame 例子、RZ spectral gridding、moving window 连续注入或 warning manager，就不会再把“参数在哪一步被锁定”与“对象在哪一步消费这些参数”混在一起。

再往前走一步，\cpath{ReadParameters()} 还不只是“把输入存进成员”。它已经在替后面的对象图做分叉。例如：

\begin{itemize}
\item \cpath{do\_fluid\_species} 先通过 \cpath{fluids.species\_names} 决定，随后直接控制是否创建 \cpath{MultiFluidContainer}；
\item \cpath{electrostatic\_solver\_id} 与 \cpath{electromagnetic\_solver\_id} 先在 \cpath{ReadParameters()} 中互相覆盖和约束，随后决定静电 solver 具体实例化成 \cpath{LabFrameExplicitES}、\cpath{EffectivePotentialES} 还是 \cpath{RelativisticExplicitES}；
\item \code{electromagnetic_solver_id == HybridPIC} 时，构造函数后半段才真正创建 \cpath{HybridPICModel}；
\item \cpath{grid\_type}、\cpath{do\_current\_centering}、\cpath{n\_rz\_azimuthal\_modes} 等参数先在 \cpath{ReadParameters()} 中完成默认值推导和合法组合检查，后面再影响 nodal flags、centering coefficients、fluid 兼容性与容器初始化。
\end{itemize}

同样，\cpath{ReadParameters()} 里还有一类“源码先推导默认，再由用户输入覆盖”的派生状态。最典型的是：

% <-- fenced code (cpp)
\begin{codeblock}
if (!do_divb_cleaning
    && m_p_ext_field_params->B_ext_grid_type != ExternalFieldType::default_zero
    && m_p_ext_field_params->B_ext_grid_type != ExternalFieldType::constant
    ...
    && (electromagnetic_solver_id == Yee
        || electromagnetic_solver_id == HybridPIC
        || ...))
{
    m_do_initial_div_cleaning = true;
}
pp_warpx.query("do_initial_div_cleaning", m_do_initial_div_cleaning);
\end{codeblock}

这里 \cpath{do\_initial\_div\_cleaning} 不是单纯从输入文件机械读取，而是先根据外部场类型、grid type 和 solver 组合推导一个默认值，再允许用户显式覆盖。也就是说，后面 \cpath{ProjectionDivCleaner} 是否参与初态构造，不是孤立地由一个布尔开关决定，而是由初始化启动层和 \cpath{ReadParameters()} 主体共同塑造的。

到这一层为止，初始化阶段已经可以更清楚地压成一条完整链：

\begin{enumerate}
\item 启动层先改 parser、锁定全局静态状态。
\item \cpath{ReadParameters()} 再把这些值落到 \cpath{WarpX} 成员，并决定对象图分叉。
\item 构造函数后半段据此创建 \cpath{MultiParticleContainer}、\cpath{ParticleBoundaryBuffer}、\cpath{MultiFluidContainer} 和 solver objects。
\item 最后 \cpath{InitData()} 才真正开始分配 level data、生成粒子和构造初态。
\end{enumerate}

但这里还差最后一层经常被误读的东西：\cpath{ReadParameters()} 里那些看上去分散的算法检查，其实会继续决定后面到底分配哪类 \cpath{MultiFab}、是否创建 implicit solver 额外状态，以及 \cpath{ProjectionDivCleaner} 是否在初态构造中插队。

最紧的一组约束是：

\begin{itemize}
\item \cpath{grid\_type=hybrid} 会把 \cpath{do\_current\_centering} 推成默认真值，并要求 \cpath{algo.field\_gathering=momentum-conserving}；
\item 这又会让 \cpath{AllocLevelData()} 中的 \cpath{aux\_is\_nodal=true}，从而把 \cpath{Efield\_aux/Bfield\_aux}、后续 coarse-aux \cpath{cax} 和 gather buffer 路径切到 nodal 版本；
\item 如果同时显式要求 \cpath{do\_current\_centering=1}，源码只允许 \cpath{grid\_type=hybrid}，并且在 level 分配时额外分配 \cpath{current\_fp\_nodal}。
\end{itemize}

换句话说，\cpath{grid\_type}、\cpath{field\_gathering\_algo}、\cpath{do\_current\_centering} 在这里并不是三个平行开关，而是一组会继续改写场 index type 和附加存储布局的上游选择器。

同样，\cpath{current\_deposition\_algo} 也不是只影响某个沉积 kernel。源码先把 \cpath{PSATD}、\cpath{HybridPIC} 和 electrostatic 的默认 deposition 拉回 \cpath{Direct}，再对 \cpath{Vay} 加上三重限制：

\begin{itemize}
\item 只能配 \cpath{PSATD}
\item 不能配 mesh refinement
\item 不能配 \cpath{do\_current\_centering}
\end{itemize}

这条限制后面会直接落成额外字段分配：如果真的选了 \cpath{Vay}，\cpath{AllocLevelMFs()} 会专门分配 \cpath{current\_fp\_vay}。所以它已经是初始化对象图的一部分，而不只是运行时算法分派。

implicit 系列更明显。\cpath{ReadParameters()} 一旦看到

\begin{itemize}
\item \cpath{semi\_implicit\_em}
\item \cpath{theta\_implicit\_em}
\item \cpath{strang\_implicit\_spectral\_em}
\end{itemize}

就会立即构造 \cpath{m\_implicit\_solver}，并同时要求：

\begin{itemize}
\item current deposition 只能是 \cpath{Esirkepov} / \cpath{Villasenor} / \cpath{Direct}
\item EM solver 只能是 \cpath{Yee} / \cpath{CKC} / \cpath{PSATD}
\item particle pusher 只能是 \cpath{Boris} / \cpath{HigueraCary}
\item field gather 不能是 momentum-conserving
\end{itemize}

然后这条链会继续穿到 \cpath{InitFromScratch()} 和 \cpath{AllocLevelMFs()}：

\begin{itemize}
\item \cpath{InitFromScratch()} 在 \cpath{mypc->InitData()} 之前调用 \cpath{m\_implicit\_solver->Define()} 和 \cpath{CreateParticleAttributes()}；
\item \cpath{AllocLevelMFs()} 在 fine patch 上额外分配 \cpath{current\_fp\_non\_suborbit} 与 \cpath{E\_old}。
\end{itemize}

因此 implicit 不是“场求解器章节内部的一个局部话题”，而是在初始化阶段就已经改变粒子属性表和字段分配合同。

\cpath{particle\_shape}、filter 和 projection div cleaning 也属于同一种“前置组合约束”。只要存在 particles 或 lasers，\cpath{algo.particle\_shape} 就变成强制参数，并直接设定 \cpath{nox=noy=noz}；这反过来继续影响 guard-cell 配额、沉积 stencil 和排序默认值。\cpath{use\_filter} 也不是后面可有可无的一步卷积，因为 \cpath{AllocLevelData()} 会先 \cpath{InitFilter()}，再把 filter stencil 长度交给 \cpath{guard\_cells.Init(...)}，于是它会继续改写 guard-cell 和 buffer 需求。

最后，\cpath{m\_do\_initial\_div\_cleaning} 也不是孤立的布尔开关。源码先根据外部 \cpath{B} 场类型、EM solver 组合和是否已经启用 \cpath{do\_divb\_cleaning} 推导默认值，再允许用户覆盖；而它的下游消费点就在初始化主链里，直接决定 \cpath{ProjectionDivCleaner} 是否参与初始场构造。所以到这一步为止，\cpath{ReadParameters()} 的真实地位可以概括成一句话：

它不是单纯读参数，而是在 \cpath{AllocLevelData()} 和 \cpath{InitFromScratch()} 之前，先把“允许哪种物理-算法-存储组合存在”这件事裁成一个可执行的对象图。

接下来再往下走一步，就会看到 \cpath{AllocLevelMFs()} 真正把这张对象图摊成一组具体字段。这里最容易误判的是 \cpath{rho}、aux、外场、HybridPIC/fluid/macroscopic/EB 这些分支。

先看 \cpath{rho\_fp}。它并不是总存在的基础字段，而是只在几类路径下显式分配：

\begin{itemize}
\item electrostatic / electromagnetostatic / effective-potential
\item \cpath{HybridPIC}
\item \cpath{do\_dive\_cleaning}
\item PSATD 且启用了 \cpath{update\_with\_rho} 或 \cpath{current\_correction}
\end{itemize}

而且它的分量数也不是固定值：普通 electrostatic 或 \cpath{HybridPIC} 只需要 \cpath{ncomps}，\cpath{do\_dive\_cleaning} 一般把它升到 \cpath{2*ncomps}，PSATD 又会根据 \cpath{JRhom} 是否开启在 \cpath{ncomps} 和 \cpath{2*ncomps} 之间切换。也就是说，WarpX 初始化阶段持不持有一份“可演化的 rho 状态”，本身就是 solver contract 的一部分。

与此平行的还有三类不同的辅助标量/约束场：

\begin{itemize}
\item \cpath{phi\_fp}：只属于 electrostatic 系列；
\item \cpath{F\_fp}：只属于 \cpath{div(E)} cleaning；
\item \cpath{G\_fp}：只属于 \cpath{div(B)} cleaning。
\end{itemize}

不要把它们混成“又分配了几个标量场”。它们分别服务于 Poisson 势解、电场散度清理和磁场散度清理，而且 \cpath{G} 的 coarse-patch index type 还会继续受 \cpath{grid\_type} 控制。

再看外场分支。源码实际上维护了两套完全不同的合同：

\begin{enumerate}
\item grid external fields
\end{enumerate}
它们分配成 \cpath{Efield\_fp\_external/Bfield\_fp\_external}，index type 必须匹配 \cpath{fp}，后面由 \cpath{AddExternalFields()} 加到主网格场上。
\begin{enumerate}
\item particle external fields
\end{enumerate}
它们分配成 \cpath{E\_external\_particle\_field/B\_external\_particle\_field}，index type 必须匹配 \cpath{aux}，而且分量数直接来自外部粒子场元数据。

所以 particle external field 不是 grid external field 的别名；它跟粒子 gather 所看的 \cpath{aux} 路径绑定得更紧。这也解释了为什么只要 \cpath{mypc->m\_E\_ext\_particle\_s} 或 \cpath{m\_B\_ext\_particle\_s} 是 \cpath{read\_from\_file}，最常见的 \code{aux -> fp} alias 优化就会被打断，\cpath{Efield\_aux/Bfield\_aux} 必须改成单独分配。

剩下几条看似“模块化”的分支，其实也都在 \cpath{AllocLevelMFs()} 里继续扩展 field registry：

\begin{itemize}
\item \code{electromagnetic_solver_id == HybridPIC} 时，\cpath{m\_hybrid\_pic\_model->AllocateLevelMFs(...)} 会追加 hybrid 自己的 level 字段；
\item \cpath{do\_fluid\_species} 时，\cpath{myfl->AllocateLevelMFs(...)} 之后立刻 \cpath{InitData(...)}，说明 fluid 已经进入初始化主链，而不是仅仅挂了个容器；
\item \code{m_em_solver_medium == Macroscopic} 时，\cpath{m\_macroscopic\_properties->AllocateLevelMFs(...)} 只允许 \cpath{lev==0}，直接把 mesh refinement 排除在外；
\item \cpath{EB::enabled()} 时，所有 level 至少都会有 \cpath{distance\_to\_eb} 与 \cpath{m\_eb\_reduce\_particle\_shape}，而 finest level 上又会继续长出 \cpath{m\_eb\_update\_E/B}；如果 solver 还是 \cpath{ECT}，则再额外长出 \cpath{edge\_lengths}、\cpath{face\_areas}、\cpath{area\_mod}、\cpath{Venl}、\cpath{ECTRhofield} 和 \cpath{FaceInfoBox} 借用关系。
\end{itemize}

因此 \cpath{AllocLevelMFs()} 的真实角色不是“机械把 \cpath{MultiFab} 开出来”，而是把前面 \cpath{ReadParameters()} 裁出来的对象图真正展开成可执行状态。

更关键的是，这些特例字段不会等到 \cpath{Evolve()} 才第一次使用。在 \cpath{InitData()} 后半段，源码会立刻：

\begin{itemize}
\item \cpath{BuildBufferMasks()}
\item \cpath{m\_macroscopic\_properties->InitData(...)}
\item \cpath{m\_electrostatic\_solver->InitData()}
\item \cpath{m\_hybrid\_pic\_model->InitData(m\_fields)}
\item \cpath{ProjectionCleanDivB()}
\end{itemize}

这就说明初始化链的最后三步其实是：

\begin{enumerate}
\item \cpath{ReadParameters()} 决定哪些组合允许存在；
\item \cpath{AllocLevelMFs()} 把这些组合摊成真实字段和对象；
\item \cpath{InitData()} 后半段立刻消费这些字段，把它们变成一份可跑的初态。
\end{enumerate}

再往后一步，\cpath{InitData()} 后半真正把这条合同闭合掉。它的顺序不是“初始化完就直接写 diagnostics”，而是先做一轮收尾和首次消费：

\begin{itemize}
\item \cpath{ComputeMaxStep()}
\item \cpath{ComputePMLFactors()}
\item 可选 \cpath{InitNCICorrector()}
\item \cpath{BuildBufferMasks()}
\item \cpath{m\_macroscopic\_properties->InitData(...)}
\item \cpath{m\_electrostatic\_solver->InitData()}
\item \cpath{m\_hybrid\_pic\_model->InitData(m\_fields)}
\item \cpath{CheckGuardCells()}
\item \cpath{ProjectionCleanDivB()}
\end{itemize}

这一步说明前面 \cpath{AllocLevelMFs()} 分配出来的对象并不是先放着，很多都会在这里立刻进入第一次初始化消费。

尤其要区分两步经常被混在一起的电场初始化动作：

\begin{enumerate}
\item \cpath{m\_electrostatic\_solver->InitData()}
\end{enumerate}
这是 solver 对象自身的初始化，fresh run 和 restart 都会走。
\begin{enumerate}
\item \cpath{ComputeSpaceChargeField(reset\_fields=false)}
\end{enumerate}
这是 fresh-run 下的初始 self-field / electrostatic solve，只在需要时触发，而且还明确保留网格上已有的用户指定值，不会先把字段清空。

它的触发条件也不是“只有 electrostatic solver 才会跑”，而是三选一：

\begin{itemize}
\item 开启 electrostatic solver
\item 任意 species 打开 \cpath{initialize\_self\_fields}
\item 指定了 boundary potential
\end{itemize}

但如果当前走的是 \cpath{HybridPIC}，源码又会显式跳过这条初始 field-solve 路径。

在这之前，若 \cpath{m\_do\_initial\_div\_cleaning} 成立，还会先执行 \cpath{ProjectionCleanDivB()}。因此初始 \cpath{B} 场的散度修正顺序其实非常明确：先完成外场读入与对象初始化，再做 \code{div B} cleaning，随后才进入初始 self-field / electrostatic solve。

Python callback 也不是一个统一“初始化结束后调用”的事件，而是被拆成了三类窗口：

\begin{itemize}
\item \cpath{beforeInitEsolve}：只在 fresh run，发生在初始场求解前；
\item \cpath{afterInitEsolve}：只在 fresh run，发生在初始场求解后、外加场提交前；
\item \cpath{afterInitatRestart}：只在 restart，表示恢复态后处理，而不是 fresh-run 初始求解窗口的一部分。
\end{itemize}

外加场这里也有两步，不能混读：

\begin{enumerate}
\item \cpath{LoadExternalFields()}
\end{enumerate}
先把 external grid fields 装到 \cpath{Efield\_fp\_external/Bfield\_fp\_external}，把 particle-only external fields 装到 \cpath{E\_external\_particle\_field/B\_external\_particle\_field}，并在 finest level 给 Python \cpath{loadExternalFields} callback 一个写这些 buffer 的机会。
\begin{enumerate}
\item \cpath{AddExternalFields()}
\end{enumerate}
再把 grid external fields 统一加回 \cpath{Efield\_fp/Bfield\_fp} 主场。

所以第 0 步最终主场不是“纯 self-field 结果”，而是“初始 self-field / electrostatic solve 结果，再叠加 grid external fields”。

这之后才进入第 0 步 diagnostics：

\begin{itemize}
\item \cpath{multi\_diags->FilterComputePackFlush(istep[0]-1)}
\item \cpath{reduced\_diags->ComputeDiags(...)}
\item \cpath{reduced\_diags->WriteToFile(...)}
\end{itemize}

因此第 0 步输出的并不是“原始输入快照”，而是“初始化主链已经全部完成后的第一份可运行状态快照”。对于 restart，这一步则取决于 \cpath{write\_diagnostics\_on\_restart}，表示是否要把恢复态也立刻写成一份 diagnostics。

\section{3A.3 顶层入口：fresh run 与 restart 分叉}

\sourceline{源码文件}{\cpath{Source/Initialization/WarpXInitData.cpp}}
\sourceline{函数}{\cpath{WarpX::InitData()}}

第 3 章已经给过总表，这里看决定数据来源的核心分支。

% <-- fenced code (cpp)
\begin{codeblock}
if (!restart_chkfile.empty())
{
    InitFromCheckpoint(restart_chkfile);
    PrintDtDxDyDz();
    PostRestart();
}
else
{
    ComputeDt();
    PrintDtDxDyDz();
    InitFromScratch();
    InitDiagnostics();
}
\end{codeblock}

这段分支决定初始化数据来源：

\begin{itemize}
\item restart：从 checkpoint 恢复 mesh、field、particles 和时间层，再做 restart 后处理；
\item fresh run：先由 CFL、网格和 solver 计算 \cpath{dt}，再从零建立 AMR level、field、particles 和 diagnostics。
\end{itemize}

物理上，restart 应该恢复一个已经离散一致的状态；fresh run 则必须从输入参数构造这种一致状态。后面讲的外部场、初始粒子、Gaussian beam、openPMD 注入和 projection cleaning，主要属于 fresh run 路径。

可以把 fresh run 与 restart 的责任边界压缩成下面的初始化路线表：

% <-- table block (source: 03a-warpx-initialization L464)
\begin{tabularx}{\textwidth}{LLL}
\toprule
输入状态 & 主要调用顺序 & 进入 \cpath{Evolve()} 前必须具备的状态 \\
\midrule
\endfirsthead
\toprule
输入状态 & 主要调用顺序 & 进入 \cpath{Evolve()} 前必须具备的状态 \\
\midrule
\endhead
checkpoint restart & \code{InitFromCheckpoint() -> PostRestart()} & 从 checkpoint 恢复 AMR level、fields、particles 和时间层，再完成 restart 后处理 \\
fresh run & \code{ComputeDt() -> InitFromScratch() -> MakeNewLevelFromScratch() -> mypc->AllocData()/InitData() -> PML -> InitDiagnostics()} & 根据输入重新建立 AMR、solver、粒子、外部场和初始 diagnostics \\
\bottomrule
\end{tabularx}

这张图的重点不是列出每一个初始化 helper，而是固定数据来源的不可互换性：restart 路径恢复已经离散化的状态，fresh run 路径才负责从参数重新物化 AMR、solver、粒子、外部场和初始 diagnostics。后面的 \cpath{PlasmaInjector}、Gaussian beam、openPMD 文件注入和 projection cleaning 都必须放在 fresh-run 分支内理解，不能被误写成 restart 的重复初始化。

\section{3A.4 \texttt{InitFromScratch()}：AMReX level 与粒子初始化}

\sourceline{源码文件}{\cpath{Source/Initialization/WarpXInitData.cpp}}
\sourceline{函数}{\cpath{WarpX::InitFromScratch()}}

AMReX 回调文件：\cpath{Source/WarpX.cpp}
回调函数：\cpath{WarpX::MakeNewLevelFromScratch()}

% <-- fenced code (cpp)
\begin{codeblock}
void
WarpX::InitFromScratch ()
{
    const Real time = 0.0;
    AmrCore::InitFromScratch(time); // calls MakeNewLevelFromScratch

    if (m_implicit_solver) {
        m_implicit_solver->Define(this, /*from_restart=*/false);
        m_implicit_solver->CreateParticleAttributes();
    }

    mypc->AllocData();
    mypc->InitData();

    InitPML();
    ExecutePythonCallback("allocdata");
}
\end{codeblock}

这里的顺序很重要：

\begin{enumerate}
\item \cpath{AmrCore::InitFromScratch(time)} 驱动 AMReX 创建 AMR level；它会回调 \cpath{WarpX::MakeNewLevelFromScratch()}，后者对每个 level 依次调用 \code{AllocLevelData(lev, ...)} 和 \code{InitLevelData(lev, time)}。因此 \cpath{AllocLevelData} 是这个回调链的一部分，不是 \cpath{InitFromScratch()} 顶层直接调用的无参步骤。
\item 若启用了 implicit solver，\cpath{Define()} 与 \cpath{CreateParticleAttributes()} 在粒子初始化前建立求解器状态和所需粒子属性。
\item \cpath{mypc->AllocData()} 为粒子容器准备数据结构，\cpath{mypc->InitData()} 创建初始粒子。
\item \cpath{InitPML()} 初始化吸收边界数据结构；随后 \cpath{allocdata} Python callback 才获得已经分配完的初始化对象。
\end{enumerate}

因此，species 初始化发生在 field/level 数据结构已经存在之后，但在正式时间推进之前。

粒子容器入口：
\sourceline{源码文件}{\cpath{Source/Particles/PhysicalParticleContainer.cpp}}
\sourceline{函数}{\cpath{PhysicalParticleContainer::InitData()}}

% <-- fenced code (cpp)
\begin{codeblock}
void PhysicalParticleContainer::InitData ()
{
    AddParticles(0); // Note - add on level 0
    Redistribute();  // We then redistribute
}
\end{codeblock}

这两行是后续所有粒子创建逻辑的入口。\cpath{AddParticles(0)} 负责生成初始粒子，\cpath{Redistribute()} 负责把粒子分配到正确 tile，并清理 invalid 粒子。

\section{3A.5 外部场初始化：grid field 与 particle external field}

WarpX 支持两类外部场：

\begin{itemize}
\item grid external field：外部场先放在 \cpath{Efield\_fp\_external/Bfield\_fp\_external} 网格 MultiFab 上，随后由 \cpath{AddExternalFields()} 叠加到初始 \cpath{E/B} 网格场并施加 field boundary；
\item particle external field：外部场保存在 \cpath{E/B\_external\_particle\_field} 中，在 particle gather 时参与粒子受力，并不经 \cpath{AddExternalFields()} 写入主 \cpath{E/B} 场。
\end{itemize}

外部场初始化的关键是：外部场可以是常量、parser 函数、openPMD 文件或 Python 回调。读者应避免把“外部场”理解成单一数组。

\sourceline{源码文件}{\cpath{Source/Initialization/WarpXInitData.cpp}}
\sourceline{函数}{\code{WarpX::LoadExternalFields(int lev)}}

% <-- fenced code (cpp)
\begin{codeblock}
if (grid_B_from_file || grid_E_from_file ||
    particle_B_from_file || particle_E_from_file) {
    WARPX_ALWAYS_ASSERT_WITH_MESSAGE(
        WarpX::do_moving_window == 0,
        "External fields from file are not compatible with the moving window.");
}
\end{codeblock}

这里的布尔名是为阅读压缩后的语义名；源码中展开检查 \cpath{B\_ext\_grid\_type/E\_ext\_grid\_type} 与 \cpath{m\_B\_ext\_particle\_s/m\_E\_ext\_particle\_s}。这个断言覆盖 \textbf{grid 和 particle 两类} file-driven external field：moving window 移动后，文件场的 MultiFab 没有自动重读/平移机制，故当前实现直接拒绝这个组合。

断言之后，\cpath{LoadExternalFields()} 分三步处理数据：parser 或文件填充 grid external field；在 finest level 调用 \cpath{loadExternalFields} Python callback；再按 metadata 的每张 map 填充 particle external field。不要把这三步压成一个无参的“读文件函数”：实际的 \code{ReadExternalFieldFromFile(path, MultiFab*, field, component, map_index)} 总是明确写出目标场与分量。

外部场的物理约束是：如果读入的是 \cpath{B} 或矢势 \cpath{A}，数值上还需要检查离散散度误差；这就是后面 projection divergence cleaner 的用途。

\section{3A.6 species 初始化：\texttt{PlasmaInjector} 是参数总容器}

\cpath{PlasmaInjector} 的职责不是推进粒子，而是把输入文件中一个 species 的初始化规则收集成一组可供 kernel 调用的对象。

\sourceline{源码文件}{\cpath{Source/Initialization/PlasmaInjector.cpp}}
\sourceline{函数}{\cpath{PlasmaInjector::PlasmaInjector()}}

% <-- fenced code (cpp)
\begin{codeblock}
std::string injection_style = "none";
utils::parser::query(pp_species, source_name, "injection_style", injection_style);
std::transform(injection_style.begin(),
               injection_style.end(),
               injection_style.begin(),
               ::tolower);

num_particles_per_cell_each_dim.assign(3, 0);

if (injection_style == "singleparticle") {
    setupSingleParticle(pp_species);
    return;
} else if (injection_style == "multipleparticles") {
    setupMultipleParticles(pp_species);
    return;
} else if (injection_style == "gaussian_beam") {
    setupGaussianBeam(pp_species);
} else if (injection_style == "nrandompercell") {
    setupNRandomPerCell(pp_species);
} else if (injection_style == "nfluxpercell") {
    setupNFluxPerCell(pp_species);
} else if (injection_style == "nuniformpercell") {
    setupNuniformPerCell(pp_species);
} else if (injection_style == "external_file") {
    setupExternalFile(pp_species);
} else if (injection_style != "none") {
    SpeciesUtils::StringParseAbortMessage("Injection style", injection_style);
}
\end{codeblock}

这个分支定义了初始粒子创建的第一层分类：

% <-- table block (source: 03a-warpx-initialization L589)
\begin{tabularx}{\textwidth}{LLL}
\toprule
\cpath{injection\_style} & 含义 & 后续创建路径 \\
\midrule
\endfirsthead
\toprule
\cpath{injection\_style} & 含义 & 后续创建路径 \\
\midrule
\endhead
\cpath{singleparticle} & 单个手工粒子 & \cpath{AddNParticles()} \\
\cpath{multipleparticles} & 手工粒子列表 & \cpath{AddNParticles()} \\
\cpath{gaussian\_beam} & 空间高斯束流 & \cpath{AddGaussianBeam()} \\
\cpath{external\_file} & openPMD 粒子文件 & \cpath{AddPlasmaFromFile()} \\
\cpath{nrandompercell} / \cpath{nuniformpercell} & 体密度注入 & \cpath{AddPlasma()} \\
\cpath{nfluxpercell} & 面通量注入 & \cpath{AddPlasmaFlux()} \\
\bottomrule
\end{tabularx}

在这一步之后，\cpath{PlasmaInjector} 还会把 host 侧 functor 拷贝到 device 侧，保证后续 GPU kernel 可以调用。

\section{3A.7 密度和动量分布：从文本参数到 functor}

密度解析由 \cpath{SpeciesUtils::parseDensity()} 完成。
\sourceline{源码文件}{\cpath{Source/Utils/SpeciesUtils.cpp}}

% <-- fenced code (cpp)
\begin{codeblock}
if (rho_prof_s == "constant") {
    amrex::Real density = 0;
    utils::parser::getWithParser(pp_species, source_name, "density", density);
    h_inj_rho.reset(new InjectorDensity(
        (InjectorDensityConstant*)nullptr, density));
} else if (rho_prof_s == "parse_density_function") {
    std::string str_density_function;
    utils::parser::Store_parserString(
        pp_species, source_name, "density_function(x,y,z)", str_density_function);
    density_parser = std::make_unique<amrex::Parser>(
        utils::parser::makeParser(str_density_function,{"x","y","z"}));
    h_inj_rho.reset(new InjectorDensity(
        (InjectorDensityParser*)nullptr, density_parser->compile<3>()));
} else if (rho_prof_s == "read_from_file") {
    std::string density_file;
    std::string field_name = "density";
    bool distributed = true;
    utils::parser::get(pp_species, source_name,
                       "read_density_from_path", density_file);
    utils::parser::query(pp_species, source_name,
                         "density_mesh_name", field_name);
    pp_species.query("read_density_distributed", distributed);
    h_inj_rho.reset(new InjectorDensity(
        (InjectorDensityFromFile*)nullptr,
        density_file, field_name, geom, distributed));
}
\end{codeblock}

因此 file profile 不只是一个路径：\cpath{density\_mesh\_name} 选择文件中的 mesh record（默认 \cpath{density}），\cpath{geom} 提供目标几何，\cpath{read\_density\_distributed} 决定读取布局。复现实验时应把这四项与文件坐标/单位一起检查，而不应只确认路径可打开。

体注入中的宏粒子权重由密度决定：

\[
w_p \approx n(\mathbf{x})\frac{\Delta V}{N_{ppc}}.
\]

动量解析由 \cpath{SpeciesUtils::parseMomentum()} 完成。常见分支包括：

\begin{itemize}
\item \cpath{at\_rest}：\cpath{u=0}；
\item \cpath{constant}：固定 \cpath{ux/uy/uz}；
\item \cpath{gaussian}：每个方向正态采样；
\item \cpath{gaussianflux}：通量注入专用，法向速度按 \code{v_n f(v)} 加权；
\item \cpath{uniform}：每个方向均匀采样；
\item \cpath{maxwell\_boltzmann}：非相对论 Maxwellian；
\item \cpath{maxwell\_juttner}：相对论热平衡分布；
\item \cpath{parse\_momentum\_function}：空间解析函数。
\end{itemize}

\cpath{InjectorMomentum} 的关键工程实现是手写 tagged union。
\sourceline{源码文件}{\cpath{Source/Initialization/InjectorMomentum.H}}
对象：\cpath{InjectorMomentum}

% <-- fenced code (cpp)
\begin{codeblock}
struct InjectorMomentum
{
    enum struct Type {
        constant,
        gaussian,
        gaussianflux,
        uniform,
        boltzmann,
        juttner,
        parser,
        gaussianparser
    };

    Type type;

    union {
        InjectorMomentumConstant constant;
        InjectorMomentumGaussian gaussian;
        InjectorMomentumGaussianFlux gaussianflux;
        InjectorMomentumUniform uniform;
        InjectorMomentumBoltzmann boltzmann;
        InjectorMomentumJuttner juttner;
        InjectorMomentumParser parser;
        InjectorMomentumGaussianParser gaussianparser;
    };
\end{codeblock}

这不是普通虚函数多态，而是 GPU kernel 友好的平铺对象。后续 \cpath{AddPlasma()} 可以在 device 上用 \cpath{getMomentum()} 采样单粒子动量，用 \cpath{getBulkMomentum()} 得到平均漂移速度。

\section{3A.8 \texttt{AddParticles()}：按注入类型进入创建函数}

\sourceline{源码文件}{\cpath{Source/Particles/ParticleCreation/AddParticles.cpp}}
\sourceline{函数}{\code{PhysicalParticleContainer::AddParticles(int lev)}}

% <-- fenced code (cpp)
\begin{codeblock}
void
PhysicalParticleContainer::AddParticles (int lev)
{
    ABLASTR_PROFILE("PhysicalParticleContainer::AddParticles()");

    for (auto const& plasma_injector : plasma_injectors) {

        if (plasma_injector->add_single_particle) {
            if (WarpX::gamma_boost > 1.) {
                MapParticletoBoostedFrame(plasma_injector->single_particle_pos[0],
                                          plasma_injector->single_particle_pos[1],
                                          plasma_injector->single_particle_pos[2],
                                          plasma_injector->single_particle_u[0],
                                          plasma_injector->single_particle_u[1],
                                          plasma_injector->single_particle_u[2]);
            }
            auto const& pos = plasma_injector->single_particle_pos;
            auto const& u = plasma_injector->single_particle_u;
            const auto weight = plasma_injector->single_particle_weight;
            const auto xp = amrex::Vector<ParticleReal>{pos[0]};
            const auto yp = amrex::Vector<ParticleReal>{pos[1]};
            const auto zp = amrex::Vector<ParticleReal>{pos[2]};
            const auto uxp = amrex::Vector<ParticleReal>{u[0]};
            const auto uyp = amrex::Vector<ParticleReal>{u[1]};
            const auto uzp = amrex::Vector<ParticleReal>{u[2]};
            const auto attr =
                amrex::Vector<amrex::Vector<ParticleReal>>{{weight}};
            const amrex::Vector<amrex::Vector<int>> attr_int;
            AddNParticles(lev, 1, xp, yp, zp, uxp, uyp, uzp,
                          1, attr, 0, attr_int, 0);
            return;
        }
\end{codeblock}

后半段分派如下：

% <-- fenced code (cpp)
\begin{codeblock}
        if (plasma_injector->gaussian_beam) {
            AddGaussianBeam(*plasma_injector);
        }

        if (plasma_injector->external_file) {
            AddPlasmaFromFile(*plasma_injector,
                              plasma_injector->q_tot,
                              plasma_injector->z_shift);
        }

        if ( plasma_injector->doInjection() ) {
            AddPlasma(*plasma_injector, lev);
        }
    }
}
\end{codeblock}

这个函数说明：\cpath{PlasmaInjector} 中可能同时保存多种初始化信息，但最终创建时按 flag 调用不同路径。

\section{3A.9 体注入 \texttt{AddPlasma()}：候选粒子、密度、动量和权重}

体注入的核心思想是：先按 cell 和 \cpath{num\_particles\_per\_cell} 创建候选粒子，然后用真实 density/bounds 筛掉无效粒子，并把有效粒子写入 SoA。

\sourceline{源码文件}{\cpath{Source/Particles/ParticleCreation/AddParticles.cpp}}
\sourceline{函数}{\cpath{PhysicalParticleContainer::AddPlasma(...)}}

% <-- fenced code (cpp)
\begin{codeblock}
// count the number of particles that each cell in overlap_box could add
amrex::Gpu::DeviceVector<amrex::Long> counts(overlap_box.numPts(), 0);
amrex::Gpu::DeviceVector<amrex::Long> offset(overlap_box.numPts());
auto *pcounts = counts.data();
amrex::Box fine_overlap_box; // default Box is NOT ok().
if (refine_injection) {
    fine_overlap_box = overlap_box & amrex::shift(fine_injection_box, -shifted);
}
amrex::ParallelFor(overlap_box, [=] AMREX_GPU_DEVICE (int i, int j, int k) noexcept
{
    const amrex::IntVect iv(AMREX_D_DECL(i, j, k));
    auto lo = getCellCoords(overlap_corner, dx, {0._rt, 0._rt, 0._rt}, iv);
    auto hi = getCellCoords(overlap_corner, dx, {1._rt, 1._rt, 1._rt}, iv);

    lo.z = applyBallisticCorrection(lo, inj_mom, gamma_boost, beta_boost, t);
    hi.z = applyBallisticCorrection(hi, inj_mom, gamma_boost, beta_boost, t);

    if (inj_pos->overlapsWith(lo, hi))
    {
        auto index = overlap_box.index(iv);
        const amrex::Long r = (fine_overlap_box.ok() && fine_overlap_box.contains(iv))?
            (AMREX_D_TERM(rrfac[0],*rrfac[1],*rrfac[2])) : (1);
        pcounts[index] = num_ppc*r;
\end{codeblock}

\cpath{applyBallisticCorrection()} 用 bulk velocity 把 boosted-frame 位置反推到 lab-frame 初始面。其公式是：

\[
z_{0,lab}=\gamma_b\left[z_{boost}(1-\beta_b\beta_z)-ct_{boost}(\beta_z-\beta_b)\right].
\]

随后用 prefix scan 得到写入 offset：

% <-- fenced code (cpp)
\begin{codeblock}
const auto max_new_particles =
    amrex::Scan::ExclusiveSum(
        counts.size(), counts.data(), offset.data());

amrex::Long pid;
{
    pid = ParticleType::NextID();
    ParticleType::NextID(pid+max_new_particles);
}
\end{codeblock}

这种“两遍法”避免在 GPU kernel 中动态 push 粒子。

真正填充粒子时，体注入权重系数为：

% <-- fenced code (cpp)
\begin{codeblock}
AMREX_GPU_HOST_DEVICE AMREX_FORCE_INLINE
amrex::Real compute_scale_fac_volume (const amrex::GpuArray<amrex::Real, AMREX_SPACEDIM>& dx,
                                      const amrex::Long pcount) {
    using namespace amrex::literals;
    return (pcount != 0) ? AMREX_D_TERM(dx[0],*dx[1],*dx[2])/pcount : 0.0_rt;
}
\end{codeblock}

即

\[
w_p = n(\mathbf{x})\frac{\Delta V}{N_{ppc}}.
\]

在 boosted-frame 分支中，密度和纵向动量做 Lorentz 变换：

% <-- fenced code (cpp)
\begin{codeblock}
dens = gamma_boost * dens * ( 1.0_rt - beta_boost*betaz_lab );
u.z = gamma_boost * ( u.z -beta_boost*gamma_lab );
\end{codeblock}

对应：

\[
n'=\gamma_b n(1-\beta_b\beta_z),\qquad
u_z'=\gamma_b(u_z-\beta_b\gamma).
\]

最后写入粒子 SoA：

% <-- fenced code (cpp)
\begin{codeblock}
u.x *= PhysConst::c;
u.y *= PhysConst::c;
u.z *= PhysConst::c;

amrex::Real weight = dens;
weight *= scale_fac;

pa[PIdx::w ][ip] = weight;
pa[PIdx::ux][ip] = u.x;
pa[PIdx::uy][ip] = u.y;
pa[PIdx::uz][ip] = u.z;
\end{codeblock}

因此 \cpath{InjectorMomentum} 返回的 \cpath{u} 是无量纲 \codeesc{\textbackslash\{\}gamma\textbackslash\{\}beta}，粒子数组中存的是 \codeesc{\textbackslash\{\}gamma v}。

\section{3A.10 Gaussian beam：显式束流列表注入}

Gaussian beam 不使用 \cpath{InjectorDensity}，而是在 IO rank 上显式随机生成粒子列表。

\sourceline{源码文件}{\cpath{Source/Particles/ParticleCreation/AddParticles.cpp}}
\sourceline{函数}{\cpath{PhysicalParticleContainer::AddGaussianBeam(...)}}

% <-- fenced code (cpp)
\begin{codeblock}
if (ParallelDescriptor::IOProcessor()) {
    // If do_symmetrize, create either 4x or 8x fewer particles, and
    // Replicate each particle either 4 times (x,y) (-x,y) (x,-y) (-x,-y)
    // or 8 times, additionally (y,x), (-y,x), (y,-x), (-y,-x)
    if (do_symmetrize){
        npart /= symmetrization_order;
    }
    // compute the weight from N_tot if the user specified npart_real = N_tot
    // compute the weight from q_tot if the user specified q_tot
    // note that npart is the number of macroparticles
    const amrex::Real weight_3d = (N_tot > 0._rt) ? (N_tot / npart) : (q_tot / (npart*charge));
\end{codeblock}

权重来自总真实粒子数或总电荷：

\[
w_{3d} =
\begin{cases}
N_{\mathrm{tot}}/N_p,
Q_{\mathrm{tot}}/(N_p q).
\end{cases}
\]

随后按高斯分布采样空间位置：

% <-- fenced code (cpp)
\begin{codeblock}
#if defined(WARPX_DIM_3D) || defined(WARPX_DIM_RZ)
    const amrex::Real weight = weight_3d;
    amrex::Real x = amrex::RandomNormal(x_m, x_rms);
    amrex::Real y = amrex::RandomNormal(y_m, y_rms);
    amrex::Real z = amrex::RandomNormal(z_m, z_rms);
#elif defined(WARPX_DIM_XZ)
    const amrex::Real weight = weight_3d/y_rms;
    amrex::Real x = amrex::RandomNormal(x_m, x_rms);
    constexpr amrex::Real y = 0._prt;
    amrex::Real z = amrex::RandomNormal(z_m, z_rms);
\end{codeblock}

如果设置 \cpath{focal\_distance}，代码用弹道近似从焦平面束斑反推出初始位置。核心公式是：

\[
t = \frac{(\mathbf{x}_f-\mathbf{x})\cdot\mathbf{n}}{\mathbf{v}\cdot\mathbf{n}},
\qquad
\mathbf{x}\leftarrow \mathbf{x}-\mathbf{v}_\perp t.
\]

同一函数在 focal-plane 分支中计算交叉时刻并反推横向初始位置：

% <-- fenced code (cpp)
\begin{codeblock}
// Compute the time at which the particle will cross the focal plane
const amrex::Real v_dot_n = v_x * n_x + v_y * n_y + v_z * n_z;
const amrex::Real t = ((x_f-x)*n_x + (y_f-y)*n_y + (z_f-z)*n_z) / v_dot_n;

// Displace particles in the direction orthogonal to the beam bulk momentum
// i.e. orthogonal to (n_x, n_y, n_z)
#if defined(WARPX_DIM_3D) || defined(WARPX_DIM_RZ)
x = x - (v_x - v_dot_n*n_x) * t;
y = y - (v_y - v_dot_n*n_y) * t;
z = z - (v_z - v_dot_n*n_z) * t;
\end{codeblock}

如果设置 symmetrization，代码为每个样本生成 4 或 8 个镜像粒子，并把权重除以阶数。这降低横向低阶统计噪声。

\section{3A.11 openPMD 粒子文件：文件粒子列表注入}

\cpath{external\_file} 路径在构造期先打开 openPMD 文件，读取可选 \cpath{charge/mass}。
\sourceline{源码文件}{\cpath{Source/Initialization/PlasmaInjector.cpp}}
\sourceline{函数}{\cpath{PlasmaInjector::setupExternalFile()}}

% <-- fenced code (cpp)
\begin{codeblock}
void PlasmaInjector::setupExternalFile (amrex::ParmParse const& pp_species)
{
#ifndef WARPX_USE_OPENPMD
    WARPX_ABORT_WITH_MESSAGE(
        "WarpX has to be compiled with USE_OPENPMD=TRUE to be able"
        " to read the external openPMD file with species data");
#endif
    external_file = true;
    std::string str_injection_file;
    utils::parser::get(pp_species, source_name, "injection_file", str_injection_file);
    // optional parameters
    utils::parser::queryWithParser(pp_species, source_name, "q_tot", q_tot);
    utils::parser::queryWithParser(pp_species, source_name, "z_shift",z_shift);
\end{codeblock}

质量和电荷优先级是：

% <-- fenced code (text)
\begin{codeblock}
input charge/mass > input species_type > openPMD charge/mass record
\end{codeblock}

真正读入粒子时，先略去只负责几何维度开关的预处理，再看每个已存在坐标的单位组合：

\begin{itemize}
\item 源码文件：\cpath{Source/Particles/ParticleCreation/AddParticles.cpp}
\item 函数：\cpath{PhysicalParticleContainer::AddPlasmaFromFile(...)}
\end{itemize}

% <-- fenced code (cpp)
\begin{codeblock}
for (auto i = decltype(npart){0}; i<npart; ++i){

    amrex::ParticleReal const weight = ptr_w.get()[i]*w_unit;

    const auto x = ptr_x.get()[i] * position_unit_x
                 + ptr_offset_x.get()[i] * position_offset_unit_x;
    const auto y = ptr_y.get()[i] * position_unit_y
                 + ptr_offset_y.get()[i] * position_offset_unit_y;
    const auto z = ptr_z.get()[i] * position_unit_z
                 + ptr_offset_z.get()[i] * position_offset_unit_z
                 + z_shift;
\end{codeblock}

openPMD 的 \cpath{position} 和 \cpath{positionOffset} 都乘以各自 \cpath{unitSI}，\cpath{z\_shift} 是 WarpX 额外偏移。

动量换算：

% <-- fenced code (cpp)
\begin{codeblock}
if (plasma_injector.insideBounds(x, y, z)) {

    // The normalized momentum is u = p / m = gamma beta c
    // with m = m_e for photons, m the particle mass otherwise.
    amrex::ParticleReal const mass_eff = (m_mass > 0.0_prt) ? m_mass : PhysConst::m_e;
    amrex::ParticleReal const ux = ptr_ux.get()[i]*momentum_unit_x/mass_eff;
    amrex::ParticleReal const uz = ptr_uz.get()[i]*momentum_unit_z/mass_eff;
    amrex::ParticleReal uy = 0.0_prt;
    if (ps["momentum"].contains("y")) {
        uy = ptr_uy.get()[i]*momentum_unit_y/mass_eff;
    }
\end{codeblock}

openPMD 文件中的 momentum 是物理动量 \cpath{p}。除以质量得到 \codeesc{p/m=\textbackslash\{\}gamma v}，这正是 WarpX 粒子数组的 \cpath{ux/uy/uz} 量纲。文件中的 \cpath{weighting} 直接成为宏粒子权重；\cpath{q\_tot} 只产生 warning，不会重标定权重。

\section{3A.12 Projection divergence cleaning：外部 \texttt{A/B} 场的初始约束修正}

如果外部加载的 \cpath{B} 或矢势 \cpath{A} 在离散网格上不满足散度约束，WarpX 可用 projection 方法清理。

数学上，给定向量场 \cpath{F}，构造：

\[
\mathbf{F}'=\mathbf{F}+\nabla_h\phi,
\qquad
\nabla_h\cdot\mathbf{F}'=0.
\]

于是

\[
\nabla_h^2\phi=-\nabla_h\cdot\mathbf{F}.
\]

\sourceline{源码文件}{}
\cpath{Source/Initialization/DivCleaner/ProjectionDivCleaner.cpp}

\sourceline{函数}{}
\cpath{ProjectionDivCleaner::setSourceFromField()}

% <-- fenced code (cpp)
\begin{codeblock}
WarpX::ComputeDivB(
    *m_source[ilev],
    0,
    {&Bx, &By, &Bz},
    WarpX::CellSize(0)
    );

m_source[ilev]->mult(-1._rt);
\end{codeblock}

然后 \cpath{ProjectionDivCleaner::solve()} 用 AMReX MLMG 解 Poisson 方程。

配置文件：
\cpath{Source/Initialization/DivCleaner/ProjectionDivCleaner.H}

% <-- fenced code (cpp)
\begin{codeblock}
amrex::MLMG mlmg(linop);
mlmg.setMaxIter(m_max_iter);
mlmg.setMaxFmgIter(m_max_fmg_iter);
mlmg.setBottomSolver(m_bottom_solver);
mlmg.setVerbose(m_verbose);
mlmg.setBottomVerbose(m_bottom_verbose);
mlmg.setConvergenceNormType(amrex::MLMGNormType::greater);
mlmg.solve({m_solution[lev].get()}, {m_source[lev].get()}, m_rtol, m_atol);
\end{codeblock}

最后修正场：

% <-- fenced code (cpp)
\begin{codeblock}
Bx_arr(i,j,k) += T::DownwardDx(sol_arr, coefs_x, n_coefs_x, i, j, k);
By_arr(i,j,k) += T::DownwardDy(sol_arr, coefs_y, n_coefs_y, i, j, k);
Bz_arr(i,j,k) += T::DownwardDz(sol_arr, coefs_z, n_coefs_z, i, j, k);
\end{codeblock}

这不是演化阶段的 \cpath{warpx.do\_dive\_cleaning/do\_divb\_cleaning}。后者是 Maxwell solver 时间推进中的清理变量或修正方程；本节讲的是初始化或外部场加载后的 Poisson projection。

\subsection{Laser antenna 与 profile 分派：laser 初始化并不走 \texttt{PlasmaInjector}}

species 初始化走的是 \code{PlasmaInjector -> AddPlasma/AddGaussianBeam/AddPlasmaFromFile} 这一条链；laser 初始化则完全不同。它的入口是：

\begin{itemize}
\item \cpath{lasers.names}
\item \cpath{LaserParticleContainer}
\item \cpath{Laser/LaserProfiles.*}
\end{itemize}

\cpath{LaserParticleContainer} 构造函数先统一读取天线几何和公共物理参数：

% <-- fenced code (cpp)
\begin{codeblock}
utils::parser::getArrWithParser(pp_laser_name, "position", m_position);
utils::parser::getArrWithParser(pp_laser_name, "direction", m_nvec);
utils::parser::getArrWithParser(pp_laser_name, "polarization", m_p_X);
utils::parser::getWithParser(pp_laser_name, "wavelength", m_wavelength);
\end{codeblock}

然后要求 \cpath{e\_max} 和 \cpath{a0} 二选一：

% <-- fenced code (cpp)
\begin{codeblock}
const bool e_max_is_specified =
    utils::parser::queryWithParser(pp_laser_name, "e_max", m_e_max);
Real a0;
const bool a0_is_specified =
    utils::parser::queryWithParser(pp_laser_name, "a0", a0);
...
AMREX_ALWAYS_ASSERT_WITH_MESSAGE(
    e_max_is_specified ^ a0_is_specified,
    "Exactly one of e_max or a0 must be specified for the laser.\n");
\end{codeblock}

如果给的是 \cpath{a0}，WarpX 会立即按

\[
E_{\max}=\frac{m_e \omega c}{q_e}a_0
\]

换算成真实场强 \cpath{e\_max}。这说明在 profile 实现层，\cpath{a0} 已经不再存在，剩下的只是规范化后的公共参数。

profile 类型分派也不是一串手写 \cpath{if-else}，而是 \cpath{LaserProfiles.H} 中的工厂字典：

% <-- fenced code (cpp)
\begin{codeblock}
laser_profiles_dictionary =
{
    {"gaussian",
        [] () {return std::make_unique<GaussianLaserProfile>();} },
    {"parse_field_function",
        [] () {return std::make_unique<FieldFunctionLaserProfile>();} },
    {"from_file",
        [] () {return std::make_unique<FromFileLaserProfile>();} }
};
\end{codeblock}

构造函数把 \cpath{profile} 字符串转成小写后直接做字典查找，再调用统一接口：

% <-- fenced code (cpp)
\begin{codeblock}
m_up_laser_profile = laser_profiles_dictionary.at(laser_type_s)();
...
m_up_laser_profile->init(pp_laser_name, common_params);
\end{codeblock}

所以：

\begin{enumerate}
\item \cpath{gaussian} 走解析包络与聚焦/STC/chirp 公式；
\item \cpath{parse\_field\_function} 直接把 \code{field_function(X,Y,t)} 编译成 parser；
\item \cpath{from\_file} 则走 \cpath{lasy\_file\_name} 或 \cpath{binary\_file\_name}，并用 \cpath{time\_chunk\_size} / \cpath{delay} 做时间分块读入。
\end{enumerate}

更重要的是，laser pulse 不是直接“写入一个初始场数组”，而是通过人工天线粒子实现。\cpath{LaserParticleContainer.H} 的类注释写得很明确：这些粒子均匀分布在一个平面上，按预设位移沉积电流 \cpath{J}，再由 Maxwell solver 在网格上生成真正的激光场。因此 \cpath{LaserParticleContainer} 需要 current deposition，但不走普通 \cpath{FieldGather}，它也正因如此直接继承 \cpath{WarpXParticleContainer}，而不是 \cpath{PhysicalParticleContainer}。

continuous injection 对 laser 还有额外几何约束：

% <-- fenced code (cpp)
\begin{codeblock}
AMREX_ALWAYS_ASSERT_WITH_MESSAGE(
    ...,
    "do_continous_injection for laser particle only works"
    " if moving window direction and laser propagation direction are the same");
\end{codeblock}

如果叠加 boosted frame，目前还进一步要求 boost 方向是 \cpath{z}。因此 laser 的 \cpath{do\_continuous\_injection} 虽然和普通 species 同名，但它的可用组合更窄，实际上是“moving-window / boosted-frame 下天线何时进入域内”的控制开关。

这还只是 laser 初始化链的入口。真正运行起来后，WarpX 并不会把解析式或文件 profile 直接写进 \cpath{E/Bfield\_fp}。它先把 profile 解释成“天线平面上每个人工粒子的目标发射振幅”，再通过标准 \cpath{J/rho} 沉积把激光交给 Maxwell solver。

\cpath{GaussianLaserProfile::init()} 继续在公共参数外读取：

% <-- fenced code (cpp)
\begin{codeblock}
utils::parser::getWithParser(ppl, "profile_waist", m_params.waist);
utils::parser::getWithParser(ppl, "profile_duration", m_params.duration);
utils::parser::getWithParser(ppl, "profile_t_peak", m_params.t_peak);
utils::parser::getWithParser(ppl, "profile_focal_distance", m_params.focal_distance);
utils::parser::queryWithParser(ppl, "zeta", m_params.zeta);
utils::parser::queryWithParser(ppl, "beta", m_params.beta);
utils::parser::queryWithParser(ppl, "phi2", m_params.phi2);
utils::parser::queryWithParser(ppl, "phi0", m_params.phi0);
\end{codeblock}

因此当前 \cpath{gaussian} profile 并不是“只有纵向和横向高斯包络”的最简形式，而是已经直接包含：

\begin{itemize}
\item 聚焦距离 \cpath{profile\_focal\_distance}
\item carrier-envelope phase \cpath{phi0}
\item spatial chirp \cpath{zeta}
\item angular dispersion \cpath{beta}
\item temporal chirp \cpath{phi2}
\end{itemize}

WarpX 随后还会把 \cpath{stc\_direction} 归一化，并强制要求它与天线法向 \cpath{nvec} 正交：

% <-- fenced code (cpp)
\begin{codeblock}
WARPX_ALWAYS_ASSERT_WITH_MESSAGE(std::abs(dp2) < 1.0e-14,
    "stc_direction is not perpendicular to the laser plane vector");
\end{codeblock}

这说明 \cpath{stc\_direction} 的真实语义是“STC 在激光平面内的作用方向”，不是随手附带的一个参考向量。到了 \cpath{fill\_amplitude()}，WarpX 再显式构造：

\begin{itemize}
\item \cpath{diffract\_factor}
\item \cpath{inv\_complex\_waist\_2}
\item \cpath{stretch\_factor}
\end{itemize}

并按 3D/RZ、XZ、1D 分别选不同 prefactor。也就是说，当前 Gaussian 实现已经把 diffraction、Gouy phase、wavefront curvature、spatial chirp、angular dispersion 和 temporal chirp 全都折叠进同一个复 envelope 公式里，而不是后面再给场求解器额外修正。

\cpath{from\_file} profile 的合同也比“读一个文件”更具体。源码强制要求：

% <-- fenced code (cpp)
\begin{codeblock}
WARPX_ALWAYS_ASSERT_WITH_MESSAGE (
    lasy_file_name.empty() != binary_file_name.empty(),
    "Exactly one of 'binary_file_name' and 'lasy_file_name' has to be specified");
\end{codeblock}

因此 \cpath{from\_file} 不是同时混用两种后端，而是在：

\begin{itemize}
\item \cpath{lasy\_file\_name}
\item \cpath{binary\_file\_name}
\end{itemize}

之间二选一。并且它不是一次性把整份文件全读进来，而是先把 \cpath{time\_chunk\_size} 设成默认全时域，再允许用户把它改小，随后只预读第一块；真正运行时由 \cpath{m\_up\_laser\_profile->update(t\_lab)} 按需换块。\cpath{delay} 也不是写文件时的预处理，而是在 \cpath{update()} / \cpath{fill\_amplitude()} 内统一平移 profile 时间轴。

更关键的是，\cpath{LaserParticleContainer::InitData()} 的产物不是“初始化场”，而是人工天线粒子。它先按 finest cell 和天线平面方向计算平面 spacing \cpath{S\_X/S\_Y}，再建立实验室坐标与激光平面坐标之间的 \cpath{Transform/InverseTransform}，最后只在注入盒覆盖的区域生成粒子：

\begin{itemize}
\item Cartesian / XZ / 1D：每个平面位置生成一对 \cpath{+w/-w} 粒子
\item RZ：围绕轴向展开成 spokes，并用 \cpath{2*pi*r/n\_spokes} 修正权重
\end{itemize}

到了 \cpath{LaserParticleContainer::Evolve()}，运行时主链是：

\begin{enumerate}
\item 若在 boosted frame，先把数值时间 \cpath{t} 转回实验室系 \cpath{t\_lab}
\item \cpath{m\_up\_laser\_profile->update(t\_lab)}，必要时推进 \cpath{from\_file} 的时间块缓存
\item \cpath{calculate\_laser\_plane\_coordinates(...)}，把真实粒子位置映回激光平面坐标
\item \cpath{fill\_amplitude(...)}，为每个天线粒子求目标 \cpath{E} 振幅
\item \cpath{update\_laser\_particle(...)}，把振幅变成粒子动量和位置
\item \cpath{DepositCurrent()} / \cpath{DepositCharge()}，把人工天线粒子写入 \cpath{current\_fp/rho\_fp}，必要时也写入 coarse-fine \cpath{current\_buf/rho\_buf}
\end{enumerate}

所以 WarpX 的 laser antenna 不是“边界直接给定场”，而是“人为构造一层发射粒子，通过普通沉积链把 profile 写成 \cpath{J}，再让 Maxwell solver 自己生成传播场”。这也解释了为什么 laser 在 mesh refinement 下照样要走 \code{fp / buf} 分流，以及为什么 \cpath{lasers.deposit\_on\_main\_grid} 其实属于 AMR/沉积合同，而不只是一个 laser 小选项。

最后，laser 的 continuous injection 也应理解得更准确一点。\cpath{ContinuousInjection()} 检查的是更新后的 \cpath{m\_updated\_position} 是否第一次进入当前 \cpath{injection\_box}；一旦进入，就调用一次 \cpath{InitData()} 生成那片天线粒子，之后再由普通 \cpath{Evolve()} 周期性更新它们的发射振幅。它不是每步“重新生成整束激光”，而是在 moving-window / boosted-frame 下决定“天线什么时候进入域内并开始工作”。

\cpath{parse\_field\_function} 这条分支也需要单独说明，因为它和前两类 profile 的数学合同并不一样。官方参数文档把它定义成：

% <-- fenced code (text)
\begin{codeblock}
<laser_name>.field_function(X,Y,t)
\end{codeblock}

这里给出的不是包络，而是完整电场本身；\cpath{X/Y} 是垂直于激光传播方向的平面坐标，而不是固定的仿真坐标轴。\cpath{FieldFunctionLaserProfile::init()} 在源码里只做两件事：

% <-- fenced code (cpp)
\begin{codeblock}
utils::parser::Store_parserString(
        ppl, "field_function(X,Y,t)", m_params.field_function);
m_parser = utils::parser::makeParser(m_params.field_function,{"X","Y","t"});
\end{codeblock}

随后 \cpath{fill\_amplitude()} 也没有再加任何 envelope、phase 或 geometry 修正，而是直接逐粒子求值：

% <-- fenced code (cpp)
\begin{codeblock}
auto parser = m_parser.compile<3>();
amrex::ParallelFor(np, [=] AMREX_GPU_DEVICE (int i) noexcept
{
    amplitude[i] = parser(Xp[i], Yp[i], t);
});
\end{codeblock}

这说明 \cpath{parse\_field\_function} 的 profile 层是三类实现里最薄的一层：用户写什么场函数，天线平面上就得到什么目标场值；后续所有“把目标场值变成可沉积粒子运动”的责任，都落在 \cpath{LaserParticleContainer} 的更新 kernel 上。

这几个 kernel 里，\cpath{calculate\_laser\_plane\_coordinates(...)} 负责把真实粒子位置减去天线参考点后，投影回 \cpath{m\_u\_X/m\_u\_Y} 基底，因此 profile 接口始终统一在激光平面坐标上。\cpath{ComputeWeightMobility()} 则先用固定的峰值速度上限 \code{eps = 0.05} 设定：

% <-- fenced code (cpp)
\begin{codeblock}
m_mobility = eps/m_e_max;
m_weight = PhysConst::epsilon_0 / m_mobility;
m_weight *= AMREX_D_TERM(1._rt, * Sx, * Sy);
\end{codeblock}

也就是说，WarpX 不是先任意给天线粒子权重，再让速度自己长出来；而是先要求峰值场下粒子速度不超过 \cpath{0.05c}，再反推单粒子权重。到了 \cpath{update\_laser\_particle()}，WarpX 再把 \cpath{amplitude[i]} 变成：

\begin{enumerate}
\item 沿主偏振方向 \cpath{p\_X} 的速度
\item boosted-frame 下额外减去沿传播方向 \cpath{nvec} 的平移速度
\item 相应的 relativistic momentum \cpath{ux/uy/uz}
\item 显式路径的整步位置推进，或 implicit 路径基于 \cpath{x\_n/y\_n/z\_n} 的半步 time-centered 位置推进
\end{enumerate}

因此 laser 的人工天线粒子并不是一个只服务显式 solver 的简单边界 hack。它同样要遵守 implicit particle-centering 合同，并继续进入普通 \cpath{DepositCurrent()/DepositCharge()} 主链。

在本书引用的 WarpX 源树中，\cpath{parse\_field\_function} 的最明确真实用例位于：

\begin{itemize}
\item 目录：\cpath{Examples/Tests/particle\_absorbing\_boundary/}
\item 输入：\cpath{inputs\_test\_1d\_particle\_absorbing\_boundary}
\end{itemize}

这个输入把：

\begin{itemize}
\item \code{laser1.profile = parse_field_function}
\item \code{laser1.field_function(X,Y,t) = ...}
\end{itemize}

嵌进了吸收边界测试里。但对应 \cpath{analysis.py} 检查的是边界附近的负向高速电子是否被抑制，而不是直接检查激光场本身。这意味着 \cpath{parse\_field\_function} 目前是“有真实 regression 入口，但没有独立 field-level 解析断言”的状态，书稿里应把这个验证边界明确写出来。

把 laser 模块整体放回 regression 版图后，还能看到另一个重要事实：不同 laser tests 的证据强度差别很大。\cpath{Examples/Tests/laser\_injection/} 的 1D/2D analysis 会直接比较 Gaussian 注入场的包络和主频；implicit 1D/2D 变体也继续复用同一组 analysis，因此并不只是“implicit 能跑通”的 checksum test。\cpath{Examples/Tests/laser\_injection\_from\_file/} 则继续给 \cpath{lasy}、legacy binary、boosted-frame 和 RZ \cpath{thetaMode} 文件提供 envelope/frequency 双断言。

但这一组还必须再分出一层 helper / prepare 边界。两个目录里的 \cpath{analysis\_default\_regression.py} 都只是 checksum helper：职责是自动识别 plotfile/openPMD 并按测试目录名调用 \cpath{evaluate\_checksum(...)}，提供历史输出基线，而不是新增 laser 物理断言。更重要的是，\cpath{laser\_injection\_from\_file/} 里那批 \cpath{inputs\_test\_*\_prepare.py} 并不是“待分析输入”，而是被 \cpath{CMakeLists.txt} 先行注册成 dependency 的外部文件生成阶段：

\begin{itemize}
\item 普通 1D/2D/3D/RZ lasy 变体统一先写 \cpath{gaussian\_laser\_3d}
\item legacy binary 变体先手工写 \cpath{gauss\_2d}
\item RZ \cpath{thetaMode} 变体先写 \cpath{laguerre\_laser\_RZ}
\end{itemize}

因此这组 regression 的正确结构不是单段输入，而是：

\begin{enumerate}
\item \cpath{prepare}
\begin{itemize}
\item 生成外部 laser 文件
\end{itemize}
\item \cpath{inject}
\begin{itemize}
\item WarpX 按 \cpath{from\_file} / \cpath{binary\_file\_name} / RZ 路径消费这些文件
\end{itemize}
\item \cpath{analysis}
\begin{itemize}
\item 再对最终包络和主频做强断言
\end{itemize}
\end{enumerate}

这条边界在解释 \cpath{Laser/} 模块时很关键，因为它说明“外部 laser 文件格式合同”本身已经是 active regression 的一部分，而不只是示例配套脚本。

\cardstep{选择激光案例：先匹配问题，再读取输入}

激光案例不应按目录数量或测试名称判断强弱。读者首先要说明自己想验证的是哪一层：

\begin{enumerate}
\item \textbf{天线是否给出了预期的场？}选 \cpath{laser\_injection/}，用 Gaussian 注入场的包络和主频作为比较量；
\item \textbf{外部 profile 文件是否被正确消费？}选 \cpath{laser\_injection\_from\_file/}，沿“生成外部文件 -> 注入 -> 包络/频率比较”三段检查；
\item \textbf{给定几何和驱动下如何搭建尾场加速？}选 \cpath{laser\_acceleration/}，先确认驱动、等离子体、moving window、边界和 diagnostics，再为场、相位或能谱另行定义 reference；
\item \textbf{输出在 boosted frame 和实验室系之间是否一致？}选 \cpath{boosted\_diags/}，比较同一 \cpath{Ez} 和粒子采样规则，而不把 writer 对齐误读成激光传播的解析验证。
\end{enumerate}

官方 \cpath{laser\_acceleration/README.rst} 把它定位为 LWFA 的 lab/boosted-frame 例子，并明确指出：接近圆柱对称时，RZ 可用较低成本捕捉相关物理；非圆激光或强 hosing 等明显非对称问题则需要 3D。于是读这些输入的第一步是选几何假设，不是从输出文件开始。

四个基础输入把这个选择具体化：

\begin{itemize}
\item 1D：moving window、连续电子注入、Gaussian laser antenna 与 \cpath{FieldProbe}；
\item 2D：PML、moving window、细化区域、连续背景电子与 Gaussian \cpath{beam}；
\item 3D：moving window、openPMD full diagnostics 与自定义粒子属性；
\item RZ：\code{n_rz_azimuthal_modes = 2}、beam/plasma 共存与 species 变量输出。
\end{itemize}

这四类设置说明如何组装不同维度下的应用问题，却不会自动给出一个共同的尾场幅度、dephasing 或能量增益结论。目录中可直接用于局部比较的脚本只有三类：

\begin{itemize}
\item \cpath{analysis\_1d\_fluid\_boosted.py}：将 1D boosted fluid WFA 的 \cpath{Ez/Jz/rho/Vz} 与理论 ODE 解比较；
\item \cpath{analysis\_refined\_injection.py}：比较细化注入时的总粒子数和 refinement edge 前方 \cpath{rho} 均匀性；
\item \cpath{analysis\_openpmd\_rz.py}：比较 RZ openPMD 的 mesh shape、species ordering 与 \cpath{rho\_<species>} 的物理中心位置。
\end{itemize}

其余多数输入采用默认输出回归。它适合发现同一输入的输出是否发生意外变化，却不能替代场幅、相位、beam loading 或激光衍射的独立比较。\cpath{README.rst} 仍将 \cpath{Analyze} 保留为 \cpath{TODO}，也正好提醒读者：应用输入提供的是搭建起点，而不是已经完成的物理说明。

\cpath{Examples/Tests/boosted\_diags/analysis.py} 对 \cpath{test\_3d\_laser\_acceleration\_btd} 检查的是 BTD plotfile 与 BTD openPMD 的 \cpath{Ez} 逐点一致，以及 \cpath{random\_fraction} 粒子子采样是否生效。它验证的是诊断重建和采样语义；不能由此推出激光包络、尾场相位或加速器性能。

因此，阅读激光相关输入时应保持一条清楚的证据链：

% <-- fenced code (text)
\begin{codeblock}
问题与几何假设
-> 选择注入、文件 profile、LWFA 或 BTD 案例
-> 明确 producer、输出和独立 reference
-> 只对该 reference 覆盖的 observable 下结论
\end{codeblock}

这样，\cpath{laser\_injection}、\cpath{from\_file}、\cpath{parse\_field\_function}、\cpath{laser\_acceleration} 与 BTD 不再是一张测试清单，而是回答不同问题的五类案例入口。

再往运行态交界看一层，laser 初始化还必须和 moving window、boosted frame、continuous injection、external fields 的更新合同一起理解。\cpath{WarpX::MoveWindow()} 在真正平移网格前，会先做三件互不等价的更新：

% <-- fenced code (cpp)
\begin{codeblock}
moving_window_x += ...;
::UpdateInjectionPosition(*mypc, gamma_boost, beta_boost, boost_direction, moving_window_dir, dt[0]);
mypc->UpdateAntennaPosition(dt[0]);
\end{codeblock}

这里：

\begin{enumerate}
\item \cpath{moving\_window\_x} 是窗口几何本身的位置；
\item \cpath{UpdateInjectionPosition(...)} 更新普通 species 的 \cpath{m\_current\_injection\_position}；
\item \cpath{UpdateAntennaPosition(dt)} 更新 laser antenna 的 \cpath{m\_updated\_position}。
\end{enumerate}

普通 species 的连续注入位置来自 \cpath{PlasmaInjector} 的 bulk momentum，再换成速度并在 boosted frame 下做洛伦兹变换；laser 则完全不走 \cpath{PlasmaInjector}，只在 \cpath{do\_continuous\_injection=1} 且 \cpath{gamma\_boost>1} 时按 boost velocity 平移天线平面。因此两者虽然都叫 continuous injection，但运行态位置更新机制不同。

两者的 \cpath{ContinuousInjection()} 语义也不同。普通物理粒子是在 moving window 新扫进来的 level-0 \cpath{particleBox} 中反复调用 \cpath{AddPlasma(...)}；laser 则只在 \cpath{m\_updated\_position} 第一次进入当前 \cpath{injection\_box} 时调用一次 \cpath{InitData()}，之后靠已有人工天线粒子在 \cpath{Evolve()} 中持续更新并沉积 \cpath{J/rho}。AMR 下这两个尺度也不同：species 的 runtime 注入盒按 level-0 cell 对齐，而 \cpath{LaserParticleContainer::InitData()} 仍然按 \cpath{maxLevel()} 的 finest spacing 建立天线粒子。

external field 的 moving-window 合同也在这里锁死。\cpath{LoadExternalFields()} 对 \code{B/E_ext_grid_type == read_from_file} 以及 particle external field 的 \cpath{read\_from\_file} 都直接断言：

% <-- fenced code (cpp)
\begin{codeblock}
WARPX_ALWAYS_ASSERT_WITH_MESSAGE(
    WarpX::do_moving_window == 0,
    "External fields from file are not compatible with the moving window." );
\end{codeblock}

原因不是 openPMD 不能读，而是 \cpath{MoveWindow()} 在新进入的 cells 上只能通过 \cpath{shiftMF(...)} 用 constant 或 parser 两种方式重建主场背景：

% <-- fenced code (cpp)
\begin{codeblock}
srcfab(i,j,k,n) = external_field;
srcfab(i,j,k,n) = field_parser(x,y,z);
\end{codeblock}

因此 constant/parser 外场可以跟着 moving window 继续生成，而 \cpath{read\_from\_file} 缺少“窗口每推进一次就按新的 physical coordinates 增量重读”的实现，所以被源码显式禁止。读者若研究文件外场，应选择固定窗口的 \cpath{load\_external\_field*} 案例；若问题是激光与 moving window 的交界，则应转向 \cpath{laser\_acceleration\_boosted}、\cpath{refined\_injection} 或 \cpath{subcycling\_mr}，再为目标 observable 定义独立比较。

\cardstep{激光在应用输入中的四种角色}

同一个 Gaussian laser 在应用输入中可能是驱动、结果量、AMR 对象或表面等离子体的起点。首先分清这个角色，才能选择正确的 observable。

\begin{enumerate}
\item \textbf{驱动固体靶：\cpath{laser\_ion}。}输入把 Gaussian laser、solid-density target、full/time-averaged diagnostics、\cpath{ParticleHistogram}、\cpath{FieldProbe} 与 \cpath{ParticleHistogram2D} 放到同一场景。\cpath{analysis\_test\_laser\_ion.py} 直接比较的是 \cpath{diagInst} 最后 5 个瞬时 \cpath{Ez} 的平均与 \cpath{diagTimeAvg} 的同一点值。因此它验证 time-averaged diagnostics 的定义；并不由此给出 TNSA cutoff energy、RPA threshold 或离子转换效率。
\item \textbf{把辐射当作结果量：\cpath{free\_electron\_laser}。}它没有 \cpath{lasers.names}，而以刚性注入电子/正电子束、boosted frame、moving window 和 undulator \cpath{B\_y(z)} 产生辐射。\cpath{analysis\_fel.py} 用 lab-frame 与 boosted-frame diagnostics 拟合 gain length，并以 FFT 比较 radiation wavelength。这里要解释的是束流辐射，而不是天线注入。
\item \textbf{检查细化区域中的天线：\cpath{laser\_on\_fine}。}该输入的关键选择是 \code{max_level = 1}、\cpath{fine\_tag\_lo/hi}、\cpath{laser1.prob\_lo/prob\_hi} 与 PML。默认输出回归只能告诉读者这些 AMR/求解器设置下输出没有意外变化；它不是传播、衍射或靶相互作用的独立物理比较。
\item \textbf{建立过密靶表面等离子体：\cpath{plasma\_mirror}。}它组合了 Gaussian laser、solid-density target、前后指数梯度、PML、field filter 和双 species 靶。该输入是研究反射率或高次谐波前的搭建起点；没有独立比较量时，不能把它称作这两类量的 benchmark。
\end{enumerate}

第二项的实现尤其容易被误读。\cpath{free\_electron\_laser} 由 \cpath{RigidInjectedParticleContainer}、\code{particles.B_ext_particle_init_style = parse_B_ext_particle_function} 与 \cpath{BackTransformed} diagnostics 组成。\cpath{zinject\_plane} 和 \cpath{rigid\_advance} 决定注入面前束团如何刚体传播；undulator 场不写入主场 \cpath{Bfield\_fp}，而是在 gather 时以 particle external field 的 \cpath{B\_y(z)} 提供。因此它的解释链是“刚性束流 -> 粒子背景场 -> BTD 恢复实验室系 -> gain length/wavelength”，不是“激光天线 -> 场注入”。

这条链的三个比较层也不能混用：\cpath{analysis\_fel.py} 比较 FEL 的 gain length 和 wavelength；\cpath{rigid\_injection} 检查刚性传播以及 plotfile/openPMD 中的束团状态；\cpath{boosted\_diags} 检查 BTD writer 对齐和 \cpath{random\_fraction} 子采样。每一层只为相应的 producer/consumer 接口提供证据。

最后，不要把 \cpath{laser\_ion} 当作“所有多物理同时开启”的总基准。它能作为 laser-target 骨架，field ionization、collisions 与 QED 可从此分叉；但这些模块分别在 \cpath{InitIonizationModule()}、\cpath{CollisionHandler} 与 \cpath{InitQED()} 接入，拥有不同的输入前提和推进时序。研究相应机制时，应到各自的 \cpath{field\_ionization/}、\cpath{collision/} 或 \cpath{qed/} 案例定义独立 observable，而不是把 time-averaged \cpath{Ez} 的比较扩大成多物理正确性证明。

\section{3A.13 初始化验证入口：哪些 regressions 真正在兜底}

前面的 3A.1-3A.12 讲的是“源码如何初始化”；若没有可执行 regression 对照，这些讲解很容易停留在静态阅读层。WarpX 对 \cpath{Initialization} 的验证并没有集中在一个目录里，而是分散在几组物理 test 中。

读这张验证地图时，始终先问四个问题：\textbf{输入创建了什么初态？比较的 observable 是什么？analysis 真正断言了什么？它没有覆盖哪条分支？} 目录名、能否运行和 checksum 都不能替代这四个问题的答案。下面按读者最常遇到的初始化任务组织入口；同一条路径若只有 checksum，便只能作为回归基线，不能推出解析正确性。

第一组是 \cpath{Langmuir}。它通常被当成 evolve 基准，但对初始化同样关键，因为它直接覆盖：

\begin{itemize}
\item \cpath{NUniformPerCell}
\item \cpath{profile=constant}
\item \cpath{parse\_momentum\_function}
\item 周期边界下的初始粒子/场一致性
\end{itemize}

例如 \cpath{analysis\_1d.py} 的主断言不是 checksum，而是把输出 \cpath{Ez} 与理论 Langmuir 波逐点比较：

% <-- fenced code (python)
\begin{codeblock}
E_sim = data[("mesh", field)].to_ndarray()[:, 0, 0]
E_th = get_theoretical_field(field, t0)
max_error = abs(E_sim - E_th).max() / abs(E_th).max()
assert error_rel < tolerance_rel
check_charge_conservation(data)
\end{codeblock}

这意味着 \cpath{Langmuir} 不只是“时间推进跑通”，而是在验证 parser 形式的初始扰动确实经过 \cpath{PlasmaInjector}、\cpath{SpeciesUtils} 和 \cpath{AddPlasma()} 正确落到了粒子和场上。

第二组是 \cpath{space\_charge\_initialization}。它最直接对应 \code{species.initialize_self_fields = 1} 这条初始化支线。输入文件显式打开：

% <-- fenced code (text)
\begin{codeblock}
beam.injection_style = "gaussian_beam"
beam.initialize_self_fields = 1
beam.momentum_distribution_type = "at_rest"
\end{codeblock}

而 analysis 脚本把输出 \cpath{Ex/Ey/Ez} 与高斯电荷团理论场直接比较。因此这组 test 实际上在硬验证：

\begin{enumerate}
\item \cpath{gaussian\_beam} 初始粒子云生成正确；
\item \cpath{InitData()} 检测到 \cpath{has\_initialize\_self\_fields}；
\item \cpath{ComputeSpaceChargeField(reset\_fields=false)} 在第一个时间步前给出了正确的 Coulomb 场。
\end{enumerate}

第三组是 \cpath{dive\_cleaning}。它验证的不是 projection cleaner，而是：

\begin{itemize}
\item 初始 Gaussian beam 状态进入演化后，
\item \code{warpx.do_dive_cleaning = 1} 与 PML 能否把 \codeesc{div(E)-rho/\textbackslash\{\}epsilon\_0} 误差传播并吸收掉，
\item 最终场是否回到理论 Gaussian beam 电场。
\end{itemize}

第四组是 \cpath{gaussian\_beam} / \cpath{external\_file}。这里要分两半看：

\begin{enumerate}
\item \cpath{analysis\_focusing\_beam.py} 和 \cpath{analysis\_rotated\_beam.py} 分别验证 \cpath{focal\_distance}、束斑统计、旋转位置和旋转动量；
\item \cpath{inputs\_test\_3d\_focusing\_gaussian\_beam\_from\_openpmd\_prepare.py} + \cpath{inputs\_test\_3d\_focusing\_gaussian\_beam\_from\_openpmd} 则覆盖 \cpath{external\_file} openPMD 粒子注入合同，包括 \cpath{weighting}、\cpath{mass}、\cpath{charge}、\cpath{positionOffset} 和 \code{momentum.unit_SI = m_e c}。
\end{enumerate}

这一组现在还能再细一层：

\begin{itemize}
\item \cpath{inputs\_test\_3d\_focusing\_gaussian\_beam\_photons} 不是新物理 benchmark，而是把同一聚焦束斑统计合同重复到 \code{species_type = photon} 路径；
\item \cpath{inputs\_test\_3d\_gaussian\_beam\_picmi.py} 则主要覆盖 PICMI \cpath{GaussianBunchDistribution} 前端到 runtime attributes 的接线，当前主要依赖 checksum，而不是独立理论断言。
\end{itemize}

这里还要诚实记录一个源码边界：\cpath{gaussian\_beam/CMakeLists.txt} 为 native \cpath{test\_3d\_focusing\_gaussian\_beam\_from\_openpmd} 指定了 \cpath{analysis.py}，但该目录没有这个文件。因而这条 native 路径只能说明 \code{prepare -> external_file inject -> diagnostics} 的输入与输出接口被覆盖，不能被称为官方的束斑物理强验证。相邻的 PICMI 入口复用 \cpath{analysis\_focusing\_beam.py}，可以用来学习应比较的束斑统计量；它也不能自动补上 native 入口缺失的 analysis。读者若修改 native 输入，应自行对照粒子数、总权重、横向均方根束斑和纵向切片统计，并预先说明容差来自何处。

第五组是 electrostatic / EB 初始化：

\begin{itemize}
\item \cpath{effective\_potential\_electrostatic} 用电子径向密度和解析 adiabatic expansion 基准比较，验证 effective-potential electrostatic solver；
\item \cpath{electrostatic\_sphere\_eb} 则用 \cpath{ChargeOnEB} reduced diag 和 \cpath{eb\_covered} 场，验证 \cpath{InitEB()}、Poisson 边界条件和带导体球的初始势问题。
\end{itemize}

最后一组是 \cpath{projection\_div\_cleaner}，它对应 3A.12 的 Poisson projection，而不是演化阶段的 \cpath{do\_dive\_cleaning}。相关 tests 已覆盖：

\begin{enumerate}
\item RZ openPMD 文件外场版本；
\item 3D PICMI 文件外场版本；
\item Python callback 版本；
\item 2D 解析外场版本。
\end{enumerate}

这些脚本的共同断言都是：初始化完成后，从 \cpath{raw} staggered \cpath{Bx\_aux/By\_aux/Bz\_aux} 重建的离散 \cpath{divB} 必须足够接近零。

这里也要区分强断言的位置：

\begin{itemize}
\item \cpath{test\_rz\_projection\_div\_cleaner} 的强断言在独立 \cpath{analysis.py} 里；
\item \cpath{test\_3d\_projection\_div\_cleaner\_picmi}、\cpath{test\_3d\_projection\_div\_cleaner\_callback\_picmi} 和 \cpath{test\_2d\_projection\_div\_cleaner\_initial\_analytical\_field\_picmi} 则都把 \cpath{divB} 断言直接写在输入脚本尾部，所以 \cpath{CMakeLists.txt} 里虽然 \cpath{analysis=OFF}，但并不等于这些条目只是 checksum-only。
\end{itemize}

再往 species 入口侧补一组，\cpath{Examples/Tests/flux\_injection/} 是直接锚定 \cpath{setupNFluxPerCell()} 的 regression 家族。这组 tests 分三条：

\begin{enumerate}
\item \cpath{analysis\_flux\_injection\_3d.py}
\begin{itemize}
\item 对 3D \cpath{NFluxPerCell} 场景同时检查总发射量、法向 Gaussian-flux 分布和切向 Gaussian 分布；
\end{itemize}
\item \cpath{analysis\_flux\_injection\_rz.py}
\begin{itemize}
\item 对 \code{flux_normal_axis = t} 的 RZ 连续注入检查粒子始终停留在预期 Larmor 半径带，并保持正确总通量；
\end{itemize}
\item \cpath{analysis\_flux\_injection\_from\_eb.py}
\begin{itemize}
\item 对 \code{inject_from_embedded_boundary = 1} 的 2D/3D/RZ 变体检查发射总数、法向/切向速度统计，以及粒子不会落入 EB 内部。
\end{itemize}
\end{enumerate}

因此 \cpath{flux\_injection} 的意义不是普通 emitter 示例，而是 \cpath{NFluxPerCell}、Gaussian-flux rejection sampling 和 embedded-boundary surface emission 这三条运行态合同的直接验证入口。

第一组是 \cpath{initial\_distribution}。它不是普通 smoke test，而是一组多 species、多分布的综合强基准：同一输入里同时覆盖

\begin{itemize}
\item \cpath{gaussian}
\item \cpath{maxwell\_boltzmann}
\item \cpath{maxwell\_juttner}
\item \cpath{gaussian\_beam}
\item parser 温度
\item parser bulk velocity
\item \cpath{uniform}
\item parser-Gaussian 动量统计
\end{itemize}

analysis 脚本把 reduced histogram、束斑统计和解析分布逐条对照。这意味着它真正验证的是 \cpath{PlasmaInjector}、\cpath{SpeciesUtils} 和 momentum-dispatch 层的 built-in / parser 初始化合同，而不只是“粒子能被建出来”。

由于初始化使用随机采样，checksum 的默认容差不应被理解为跨机器、跨 MPI 布局都逐位相同的物理合同。复现实验应优先比较 analysis 脚本定义的直方图、均值、方差、束斑等统计量，并在报告中写清样本量、随机数种子（若可控）、并行布局和容差；只有这些条件一致时，checksum 才能作为补充证据。

\subsection{3A.13.1 初始化验证卡：分布统计和初始自场是两份合同}

初态常被一句“束流已经初始化”概括，但这至少混合了两件不同的事：粒子的位置、权重和动量是否来自预期分布；以及这些粒子是否已经在第一步前建立了应有的自场。前者即使完全正确，也不推出 Poisson 或 space-charge 初始化正确；后者即使场形状正确，也不能替代对粒子分布和总电荷的检查。读者应把它们拆成两条 producer -> observable -> consumer 链。

\textbf{合同 A：粒子分布。}\cpath{initial\_distribution} 的输入设定 \code{max_step = 0}，并以 \cpath{ParticleHistogram} reduced diagnostics 输出各 species 的动量或位置直方图。例如 \cpath{h1x/h1y/h1z} 对 Gaussian species 统计 \cpath{ux/uy/uz}，\cpath{h4x/h4y/h4z} 对 Gaussian beam 统计空间坐标；同一组还写出 beam 的总电荷。CTest 将 producer 接到 \cpath{analysis.py}，该 consumer 读取这些 reduced diagnostics，将 Gaussian、Maxwellian、Maxwell-Juttner 和 beam 的位置统计分别与解析密度比较，并对归一化误差及相对电荷误差断言容差。

这条设计的关键价值在于 \code{max_step = 0}：它把比较对象限定为初始化完成后、尚未由时间推进改变的粒子统计。因此它能支持的结论是“给定输入、采样数、诊断定义和分析容差下，粒子初始化分布与指定解析分布相符”。它不能支持“初始空间电荷场已经正确”或“随后任意步数、任意求解器和任意 MPI 布局的动力学都正确”。若读者修改的是 \cpath{momentum\_distribution\_type}、parser mean/std、束斑 rms、截断或粒子数，应先复用这一类 histogram/charge observable；不要只看 plotfile 是否写出。

\textbf{合同 B：初始自场。}\cpath{space\_charge\_initialization} 选择静止 Gaussian beam，并显式给出：

% <-- fenced code (text)
\begin{codeblock}
beam.injection_style = "gaussian_beam"
beam.initialize_self_fields = 1
beam.momentum_distribution_type = "at_rest"
\end{codeblock}

这时 producer 不止是粒子生成。\cpath{WarpX::InitData()} 的 fresh-run 分支会遍历 species；只要发现任一 \cpath{initialize\_self\_fields} 为真，便在第一步推进之前调用 \code{ComputeSpaceChargeField(reset_E_field=false, reset_B_field=false)}。输入还请求 Full diagnostics 写出 \cpath{Ex/Ey/Ez}，CTest 将 \cpath{diags/diag1000001} 交给 \cpath{analysis.py}。这个 consumer 按二维或三维 Gaussian charge distribution 构造理论场，再对输出的 \cpath{Ex/Ey}，以及三维时的 \cpath{Ez} 做逐数组 \cpath{allclose} 容差检查。

因此这条合同能支持的结论是“在该静止 Gaussian-beam、边界、网格、shape、两 rank 注册布局和字段比较容差下，初始化自场与该理论场一致”。它不把 \cpath{initial\_distribution} 的 histogram 通过自动升级为场求解验证，也不覆盖相对论自场、open-boundary FFT、embedded boundary 或后续 pusher 的正确性；这些问题应分别进入 \cpath{relativistic\_space\_charge\_initialization}、\cpath{open\_bc\_poisson\_solver}、\cpath{magnetostatic\_eb} 或 \cpath{repelling\_particles} 的 consumer。

\textbf{实际操作顺序。}当一个新输入同时改了束流分布和 \cpath{initialize\_self\_fields} 时，先在零步或等价的初态输出上锁定合同 A，再在明确的初始 field diagnostic 上锁定合同 B；随后才用轨迹、能量或束斑演化回答“该初场被 pusher 消费后会怎样”。若 A 失败，先查 injection style、parser、随机采样、权重和 histogram 定义；若 A 通过而 B 失败，转查边界、solver、field diagnostic、\cpath{initialize\_self\_fields} 与初始场求解。两类失败的定位入口不同，不能合并成“初始化失败”。

第二组是 \cpath{initial\_plasma\_profile}。这组当前没有独立 \cpath{analysis.py}，只有 checksum helper，但输入本身非常明确：

\begin{itemize}
\item \code{injection_style = NUniformPerCell}
\item \code{profile = parse_density_function}
\end{itemize}

并把横向 parabolic channel 与纵向 ramp / plateau / ramp 组合成二维电子密度。所以它更准确地是：

\begin{itemize}
\item \cpath{parse\_density\_function} 抛物型通道初始化的 checksum-only 基线
\end{itemize}

而不是应继续留在 \code{general / to classify} 的未知条目。

再往 \cpath{initialize\_self\_fields} 这一支补一组，\cpath{repelling\_particles} 是一个更小但更干净的两体基准。它只放两个同号 \cpath{SingleParticle} species，却同时打开：

\begin{itemize}
\item \code{electron1.initialize_self_fields = 1}
\item \code{electron2.initialize_self_fields = 1}
\end{itemize}

analysis 随后从连续 plotfiles 读取两粒子的间距和速度，并用两体排斥的非相对论能量守恒关系构造理论 \codeesc{\textbackslash\{\}beta(d)}。因此这组 regression 的真实意义不是“一对粒子大概会分开”，而是：

\begin{enumerate}
\item 初始 electrostatic self-field 确实被建立起来；
\item 后续 pusher 对这份初始场的消费是对的；
\item 两者联立后能回到解析两体减速关系。
\end{enumerate}

接着把另外三组容易落单的入口也补上之后，初始化验证图还能再细一层。

第一组是 \cpath{load\_external\_field}。它不是普通粒子轨道测试，而是在验证两套初始化合同：

\begin{enumerate}
\item grid external field：
\begin{itemize}
\item \cpath{LoadExternalFields()}
\item \cpath{ReadExternalFieldFromFile()}
\item \cpath{Bfield\_fp\_external/Efield\_fp\_external}
\item \cpath{AddExternalFields()}
\end{itemize}
\item particle external field：
\begin{itemize}
\item \cpath{Particles/ExternalParticleFields.cpp}
\item \code{m_B_ext_particle_s / m_E_ext_particle_s}
\item \cpath{B\_external\_particle\_field/E\_external\_particle\_field}
\item \cpath{GetExternalFields.cpp}
\end{itemize}
\end{enumerate}

\cpath{analysis\_3d.py} / \cpath{analysis\_rz.py} 通过磁镜中的单粒子最终位置做硬断言，说明这组 regression 同时验证了“初始化写场”和“后续 gather 消费场”的接口契合。时间依赖变体 \cpath{analysis\_time\_scaling.py} 则不看粒子，而是直接比较两个时刻 plotfile 上 \cpath{B} 分量的缩放比，从而验证 \cpath{read\_fields\_*\_dependency(t)} parser 和多 field map 的时间缩放合同。

这组 family 里还要再区分一层 restart 保真。当前活跃的三条最小入口是：

\begin{itemize}
\item \cpath{test\_3d\_load\_external\_field\_particle\_time\_restart}
\item \cpath{test\_rz\_load\_external\_field\_grid\_restart}
\item \cpath{test\_rz\_load\_external\_field\_particles\_restart}
\end{itemize}

它们都只是：

\begin{itemize}
\item 先继承对应非 restart 输入
\item 再通过 \code{amr.restart = ../.../chk000150} 从中间 checkpoint 恢复
\item 然后复用 \cpath{analysis\_default\_restart.py} 逐字段比较 restart 与非 restart 输出
\end{itemize}

因此这三条 regression 验证的不是新的外场 physics，而是：

\begin{enumerate}
\item \cpath{read\_from\_file} 的 grid external field 状态能否在 restart 后保持一致；
\item \cpath{read\_from\_file} / dependency parser 的 particle external field 状态能否在 restart 后保持一致；
\item 初始化阶段构造出的 external-field 寄存器，不会在 checkpoint/restart 边界上丢失或漂移。
\end{enumerate}

第二组是 \cpath{relativistic\_space\_charge\_initialization}。它的输入和普通 \cpath{space\_charge\_initialization} 一样也打开了：

% <-- fenced code (text)
\begin{codeblock}
beam.initialize_self_fields = 1
beam.injection_style = "gaussian_beam"
\end{codeblock}

但束流动量改成了 relativistic：

% <-- fenced code (text)
\begin{codeblock}
beam.uz = 100.0
\end{codeblock}

因此这组 regression 实际验证的是 \cpath{RelativisticExplicitES::ComputeSpaceChargeField()}，而不再只是静止高斯电荷团的 Coulomb 场。analysis 脚本把 \cpath{Ex} 与理论值比较，并检查 \cpath{By} 是否满足相对论束流自场的 \codeesc{By \textbackslash\{\}approx Ex/c} 结构。这说明 \cpath{initialize\_self\_fields} 在 relativistic solver 分支下对应的是另一份初始化合同。

第三组是 \cpath{open\_bc\_poisson\_solver}。它把四个条件绑在一起：

% <-- fenced code (text)
\begin{codeblock}
boundary.field_lo = open open open
boundary.field_hi = open open open
warpx.do_electrostatic = relativistic
warpx.poisson_solver = fft
electron.initialize_self_fields = 1
\end{codeblock}

同时粒子不是 \cpath{gaussian\_beam}，而是 \code{parse_density_function + NUniformPerCell}。analysis 脚本用 Basseti-Erskine 公式逐个 \cpath{z} 截面比较 \cpath{Ex/Ey}，因此它验证的不是一般的 electrostatic 初始化，而是：

\begin{itemize}
\item open boundary
\item relativistic bunch
\item FFT Poisson
\item 以及可选 \code{warpx.use_2d_slices_fft_solver = 1}
\end{itemize}

共同定义的初始 Poisson 解是否正确。

接下来三组补足了文件驱动分布、带嵌入边界的自场，以及 collocated 采样的分支；它们与前面的束流注入和 Poisson 例子不共享同一个可观测量，因而应分别阅读。

第一层是 \cpath{load\_density}。这组 regression 的输入明确使用：

% <-- fenced code (text)
\begin{codeblock}
electrons.profile = "read_from_file"
electrons.read_density_from_path = "../test_*_load_density_prepare/example-density.h5"
electrons.do_continuous_injection = 1
warpx.do_moving_window = 1
\end{codeblock}

源码上它直接对应 \cpath{SpeciesUtils::parseDensity()} 里 \code{profile = read_from_file} 分支，把 openPMD density mesh 装进 \cpath{InjectorDensityFromFile}，然后交给 \cpath{PlasmaInjector} 和连续注入主链消费。analysis 脚本则逐个 iteration 读取 diagnostics 中的 \cpath{rho}，与 prepare 脚本定义的 ramp / parabolic channel profile 比较。因此 \cpath{load\_density} 验证的不是一般 I/O，而是“file-driven density profile + moving-window 连续注入”这条初始化合同。

第二层是 \cpath{magnetostatic\_eb}。原生 inputs 文件把：

% <-- fenced code (text)
\begin{codeblock}
warpx.do_electrostatic = labframe-electromagnetostatic
beam.initialize_self_fields = 1
warpx.eb_implicit_function = "(x**2+y**2-radius**2)"
warpx.eb_potential(x,y,z,t) = "1."
\end{codeblock}

绑在一起，而 \cpath{WarpXInitData.cpp} 的 fresh-run 分支会在 \cpath{ComputeSpaceChargeField(reset\_fields)} 之后继续调用 \cpath{ComputeMagnetostaticField()}。所以这组 test 验证的是 embedded boundary、边界 potential、初始 self-field 和 magnetostatic solve 在初始化阶段的联动。这里还必须区分两层证据：原生 \cpath{inputs\_test\_3d\_magnetostatic\_eb} 目前主要由 checksum 兜底；两个 PICMI 输入文件则在 \cpath{sim.step()} 之后内嵌了解析 \codeesc{E\_r/B\_\textbackslash\{\}theta} 误差断言，因此它们不是简单的 checksum-only regression。

第三层是 \cpath{nodal\_electrostatic}。这组输入把：

% <-- fenced code (text)
\begin{codeblock}
warpx.do_electrostatic = relativistic
warpx.grid_type = collocated
beam_p.initialize_self_fields = 1
beam_p.do_qed_quantum_sync = 1
\end{codeblock}

放在同一条链上。它的 analysis 不直接比较 \cpath{E/B}，而是用 reduced diagnostics 断言 \cpath{ParticleExtrema\_beam\_p} 给出的最大 \cpath{chi} 极小，且 \cpath{ParticleNumber} 中 photon 数始终为零。也就是说，这组 regression 验证的是 collocated relativistic electrostatic 初始 self-field 没有制造出会假触发 QED 的非物理场，它更准确地是一个“零触发基准”。

这样 \cpath{nodal\_electrostatic}、\cpath{open\_bc\_poisson\_solver} 和 \cpath{relativistic\_space\_charge\_initialization} 就不再需要继续共用一个过粗的 \code{electrostatic / Poisson} 桶。它们分别对应的是：

\begin{itemize}
\item collocated relativistic electrostatic 零触发基准
\item open boundary + FFT/sliced FFT 的 relativistic Poisson 初始化
\item relativistic Gaussian beam 的初始 self-field
\end{itemize}

不过还有一组此前仍容易被写得过粗：\cpath{electrostatic\_sphere}。它不该和一般 \cpath{Poisson} 条目混成一桶，因为它真正验证的是“一个静止均匀电子球在自身 Coulomb 场下膨胀”这条自场初始化主链。\cpath{analysis\_electrostatic\_sphere.py} 的第一层断言是解析电场对照：脚本用最终输出时间 \cpath{t\_max} 反解球半径 \cpath{r\_end}，再构造内外球解析 \cpath{E(r)}，沿坐标轴比较 \cpath{Ex/Ey/Ez} 或 RZ 下的 \cpath{Er/Ez} 的相对 \cpath{L2} 误差。这意味着它验证的不只是 solver 跑通，而是：

\begin{enumerate}
\item 初始电子球几何是否被正确装填；
\item 初始自场是否正确建立；
\item 后续 electrostatic 演化是否仍与解析膨胀解一致。
\end{enumerate}

这组 test 的第二层断言只在 lab-frame 变体上才打开。只有当输入显式要求：

% <-- fenced code (text)
\begin{codeblock}
warpx.do_electrostatic = labframe
diag2.electron.variables = x y z ux uy uz w phi
\end{codeblock}

analysis 才会利用粒子 \cpath{phi} 重建：

\begin{itemize}
\item 动能
\item 自场势能 \codeesc{0.5 \textbackslash\{\}sum w q \textbackslash\{\}phi}
\end{itemize}

并检查初末总能量。也就是说：

\begin{itemize}
\item 所有 \cpath{electrostatic\_sphere} 变体都做解析电场对照；
\item 只有写出 \cpath{phi} 的 lab-frame 版本额外验证能量账本。
\end{itemize}

各输入变体的角色也不同：

\begin{itemize}
\item \cpath{inputs\_test\_3d\_electrostatic\_sphere}
\begin{itemize}
\item 3D relativistic electrostatic 自场膨胀基线
\end{itemize}
\item \cpath{inputs\_test\_3d\_electrostatic\_sphere\_lab\_frame}
\begin{itemize}
\item lab-frame 自场膨胀 + 能量守恒
\end{itemize}
\item \cpath{inputs\_test\_3d\_electrostatic\_sphere\_lab\_frame\_mr\_emass\_10}
\begin{itemize}
\item lab-frame + MR；\code{electron.mass = 10} 主要是减慢膨胀，便于短步数比较
\end{itemize}
\item \cpath{inputs\_test\_3d\_electrostatic\_sphere\_rel\_nodal}
\begin{itemize}
\item \code{warpx.grid_type = collocated}，验证 collocated electrostatic 布局
\end{itemize}
\item \cpath{inputs\_test\_3d\_electrostatic\_sphere\_adaptive}
\begin{itemize}
\item adaptive-dt 变体；analysis 仍用实际 \cpath{t\_max} 做强对照
\end{itemize}
\item \cpath{inputs\_test\_rz\_electrostatic\_sphere}
\begin{itemize}
\item RZ 自场膨胀 + 能量守恒
\end{itemize}
\item \cpath{inputs\_test\_rz\_electrostatic\_sphere\_uniform\_weighting}
\begin{itemize}
\item 再额外覆盖 \code{radial_numpercell_power = 1.} 的 uniform-weighting 粒子装填
\end{itemize}
\end{itemize}

另外，这组目录下还有两个容易误读的文件：

\begin{itemize}
\item \cpath{analysis\_default\_regression.py} 只是通用 checksum helper；
\item \cpath{catalyst\_pipeline.py} 只是 ParaView Catalyst 的可视化脚本。
\end{itemize}

因此，\cpath{electrostatic\_sphere} 在初始化验证图里更准确的位置应单独写成：

\begin{enumerate}
\item electrostatic self-field expansion：\cpath{electrostatic\_sphere*}
\end{enumerate}

还剩另外两组此前也容易被写得过粗，但它们和 \cpath{electrostatic\_sphere} 不是一类问题。

第一组是 \cpath{electrostatic\_dirichlet\_bc}。这组不是带电粒子自场 benchmark，而是纯粹在验证时变边界势是否真正进入 electrostatic 解。输入直接设置：

% <-- fenced code (text)
\begin{codeblock}
warpx.do_electrostatic = labframe
boundary.potential_lo_x = 150.0*sin(2*pi*6.78e+06*t)
boundary.potential_hi_x = 450.0*sin(2*pi*13.56e+06*t)
diag1.fields_to_plot = phi
\end{codeblock}

analysis 并不读取内部粒子量，而是逐个输出时刻取两侧边界上的平均 \cpath{phi}，然后与理论正弦函数比较。因此它测的不是一般 \cpath{Poisson} 正确性，而是：

\begin{itemize}
\item \code{boundary.potential_lo_x / potential_hi_x}
\item time-dependent parser
\item electrostatic Dirichlet boundary condition
\item diagnostics 中 \cpath{phi} 的边界保真
\end{itemize}

PICMI 变体只是在 \code{Cartesian2DGrid(..., lower_boundary_conditions=["dirichlet", ...])}、\cpath{warpx\_potential\_lo\_x}、\cpath{warpx\_potential\_hi\_x} 这一层重复同一件事，所以它更准确的归类是：

\begin{enumerate}
\item time-dependent electrostatic boundary driving：\cpath{electrostatic\_dirichlet\_bc*}
\end{enumerate}

第二组是 \cpath{effective\_potential\_electrostatic}。这组当前只有一个 PICMI test，而且它验证的也不是一般意义上的 electrostatic \cpath{phi} 场，而是 \cpath{warpx\_effective\_potential=True} 这条 solver 分叉是否能在导体球约束下复现绝热膨胀 benchmark。PICMI 输入脚本同时做了三件关键事：

\begin{enumerate}
\item 用 \cpath{GaussianBunchDistribution} 初始化电子和离子球团；
\item 加入导体球 embedded boundary；
\item 选择：
\end{enumerate}

% <-- fenced code (python)
\begin{codeblock}
picmi.ElectrostaticSolver(
    method="Multigrid",
    warpx_effective_potential=True,
    warpx_effective_potential_factor=C_EP,
)
\end{codeblock}

analysis 则从 \cpath{sim\_parameters.dpkl} 读回参数，构造 Connor et al. 风格的绝热膨胀近似电子密度，再把 openPMD 中的 \cpath{rho\_electrons} 变成径向密度曲线，与理论值逐个输出时刻比较 RMS 误差。因此它的真实验证对象是：

\begin{itemize}
\item effective-potential electrostatic solver
\item PICMI front-end
\item 导体球约束下的电子密度膨胀近似
\end{itemize}

更准确的归类应是：

\begin{enumerate}
\item effective-potential electrostatic：\cpath{effective\_potential\_electrostatic}
\end{enumerate}

\section{3A.14 从 Birdsall \texttt{3A ES1} 到 WarpX：历史最小程序骨架的现代映射}

Birdsall and Langdon 的 \code{3A ES1} 是一份很好的历史参照，因为它把一维静电 PIC 的最小程序压缩成一条容易检查的阶段链：

% <-- fenced code (text)
\begin{codeblock}
INIT -> SETRHO -> FIELDS -> SETV -> ACCEL -> MOVE -> HISTRY
\end{codeblock}

这条链可以帮助读者理解“初始化结束后第一步推进究竟从哪里开始”，但不能把旧程序的子程序名直接当成 WarpX 的函数名。历史参照见 Birdsall 与 Langdon 的 \emph{Plasma Physics via Computer Simulation}；现代实现应从 \cpath{Source/Initialization/WarpXInitData.cpp}、\cpath{Source/Particles/} 和 \cpath{Source/FieldSolver/} 的对应职责回查。

% <-- table block (source: 03a-warpx-initialization L1727)
\begin{tabularx}{\textwidth}{LL}
\toprule
\code{3A ES1} 阶段 & WarpX 的现代映射与不能直接等同的部分 \\
\midrule
\endfirsthead
\toprule
\code{3A ES1} 阶段 & WarpX 的现代映射与不能直接等同的部分 \\
\midrule
\endhead
\cpath{INIT} & 建立网格、粒子、权重、边界和初始参数，对应 \cpath{ReadParameters()}、\cpath{WarpX::InitData()}、\cpath{InitFromScratch()}、\cpath{AllocLevelData()} 与 \cpath{mypc->AllocData()}；WarpX 还要处理 AMR、多几何、solver 分支、PML、外部场、restart 和并行布局。 \\
\cpath{SETRHO} & 用初始位置形成源，对应 \cpath{PlasmaInjector}/\cpath{AddParticles} 后的初始 \cpath{rho} 构造及 field-register 路径；\cpath{rho} 是否直接写入、重新沉积或被 solver 消费，取决于 geometry、solver 和初始化选项。 \\
\cpath{FIELDS} & 从 \cpath{rho} 求势和场，对应 electrostatic solver 的 \cpath{InitData()}、\cpath{ComputeSpaceChargeField()}、初始场填充和 projection cleaning；不能简化成单一 FFT Poisson 路径，EM、PSATD、RZ、EB 与 external field 会改变对象和约束。 \\
\cpath{SETV} & 设置热分布和漂移，对应 \cpath{SpeciesUtils}、\cpath{InjectorMomentum}、temperature/velocity functor 与粒子属性创建；WarpX 还可能创建 relativistic、spin、implicit 或 pusher-specific attributes。 \\
\cpath{ACCEL} & 用场更新速度，对应 \cpath{Evolve()} 内 particle push 与 gather；历史静电推进不能覆盖现代 EM、Boris/Vay/Higuera-Cary、implicit 和 subcycling。 \\
\cpath{MOVE} & 用新速度更新位置，对应 \cpath{PushParticlesAndDeposit()} 的 position update、边界处理和 current deposition；WarpX 还要处理 AMR tile、moving window、particle boundary、suborbit/crossing 和 MPI 交换。 \\
\cpath{HISTRY} & 记录历史量，对应 full/reduced diagnostics、openPMD/plotfile writer 与 reader-side analysis；现代 diagnostics 是独立 writer/schema 合同，不是旧程序中的一个历史数组。 \\
\bottomrule
\end{tabularx}

最容易误读的是 \code{SETRHO -> FIELDS -> SETV} 的顺序。对历史 ES1 来说，它表达的是“先由初始位置形成源，再求场，再给粒子速度”；对 WarpX fresh run，实际初始化链还要先完成 AMReX level 和 field data 分配，并根据 solver、外场和 \cpath{initialize\_self\_fields} 等条件决定哪些源/场对象被创建。因而更可靠的读法是：Birdsall 骨架提供物理阶段的语义，WarpX 源码决定这些语义在现代对象图和分支条件中的落点。

这条映射也解释了为什么 \cpath{Langmuir} regression 可以同时作为初始化和演化入口：它检查的不是某个名为 \cpath{SETRHO} 的函数，而是初始粒子/密度、初始场、后续推进和最终解析频率之间的组合合同。相反，\cpath{initial\_distribution}、\cpath{space\_charge\_initialization}、\cpath{load\_external\_field} 与 \cpath{projection\_div\_cleaner} 分别覆盖粒子分布、初始 self-field、外部场装填和初始散度修正的更窄路径。它们共同支撑的是现代 WarpX 初始化链的分层验证，而不是对 \code{3A ES1} 原程序做逐行复现。

\section{3A.15 本章小结：初始化状态怎样进入第一步推进}

到 \cpath{InitData()} 结束时，WarpX 已经完成：

% <-- fenced code (text)
\begin{codeblock}
inputs
  -> WarpX::ReadParameters()
  -> WarpX::InitData()
  -> InitFromScratch or InitFromCheckpoint
  -> AMReX levels and WarpX field data
  -> external fields
  -> species PlasmaInjector
  -> AddParticles / AddPlasma / AddGaussianBeam / AddPlasmaFromFile
  -> projection div cleaner if needed
  -> initial diagnostics
  -> Evolve()
\end{codeblock}

这个状态才是 PIC 时间推进的初始条件。后续 \cpath{Evolve()}、\cpath{OneStep()}、\cpath{PushParticlesandDeposit()}、field solver 和 diagnostics 都在这个已经离散化、分布式、带权重和边界条件的状态上工作。

初始化到推进的交界由 \cpath{Evolve()}、\cpath{OneStep()} 与 \cpath{PushParticlesandDeposit()} 承接：前者负责外层时间步，第二者负责 solver/AMR 分派，最后才把已经初始化的粒子和场送入实际粒子推进与沉积路径。历史 ES1 的阶段名只提供物理职责的参照，不能替代这些现代对象和分支。

进一步学习可沿以下方向展开：

\begin{itemize}
\item 把 \cpath{InitLevelData()} 中每一类 field allocation 展开到 root/fieldsolver 章节；
\item 把 Gaussian beam 的 emittance/focal distance 公式结合 accelerator beam optics 文献继续推导；
\item 把 openPMD 文件格式与 WarpX 单位约定加入诊断/I/O 章节；
\item 为不同初始化路径设计可重复的验证案例，并记录它们各自能说明的边界。
\end{itemize}

\section{3A.16 练习与最小复现}

\begin{enumerate}
\item \textbf{fresh/restart 定位题}：沿 \cpath{WarpXInitData.cpp} 追踪 \cpath{InitData()}，说明 \cpath{ComputeDt()} 为什么只出现在 fresh-run 分支，以及 restart 为什么必须进入 \cpath{PostRestart()}。
\item \textbf{初始化与复现题}：解释 \cpath{AmrCore::InitFromScratch()}、\cpath{AllocLevelData()}、\cpath{mypc->AllocData()}、\cpath{mypc->InitData()} 和 \cpath{InitPML()} 的先后关系。再在 \cpath{Examples/Tests/initial\_distribution/} 选择匹配输入，核对生效输入、首个 diagnostics 和 species 数，并说明输入不匹配导致的 abort 为什么既不构成初始化通过，也不构成其反证。
\end{enumerate}
